\documentclass[11pt,a4paper,oneside]{book}

\usepackage{heldenpfade}

\newcommand{\HPDocumentVersion}{unveröffentlicht}
\newcommand{\HPDocumentDate}{ohne Veröffentlichungsdatum}
\newcommand{\HPDocumentStatus}{Version}
\IfFileExists{shared/version-info.generated.tex}{%
  \input{shared/version-info.generated.tex}%
}{}

\title{Heldenpfade}
\author{}
\date{}

\begin{document}

\frontmatter
\begin{titlepage}
  \centering
  \vspace*{0.22\textheight}
  {\Huge\bfseries Heldenpfade\par}
  \vfill
  {\small
    \HPDocumentStatus{} \HPDocumentVersion\par
    \vspace{0.35em}
    \HPDocumentDate\par
  }
\end{titlepage}
\tableofcontents

\mainmatter
\chapter{Was ist Heldenpfade?}

\section*{Hinweis zur Sprache}

Zur besseren Lesbarkeit wird in diesem Regelwerk überwiegend die Bezeichnung „Spieler“ verwendet. Gemeint sind dabei selbstverständlich immer alle Spielerinnen und Spieler. Gleiches gilt für Begriffe wie „Held“ oder „Spielleitung“, die unabhängig vom Geschlecht verwendet werden.

\section*{Die wichtigste Regel}

Heldenpfade soll allen am Spieltisch Spaß machen.

Die Regeln dieses Regelwerks bilden einen gemeinsamen Rahmen. Sie sollen dabei helfen, spannende Abenteuer zu erleben und gemeinsam zu entscheiden, wie die Geschichte weitergeht.

Wenn ihr beim Spielen merkt, dass eine andere Lösung oder Regel eurer Gruppe mehr Spaß macht, dann ändert diese Regel einfach. Die Spielleitung entscheidet gemeinsam mit den Spielern, welche Lösung am besten zu eurer Geschichte und eurer Gruppe passt.

Es geht bei Heldenpfade nicht darum zu gewinnen oder zu verlieren. Das Ziel ist es, gemeinsam eine spannende Geschichte zu erleben.

In Heldenpfade erlebt ihr gemeinsam Abenteuer.

Ihr erkundet geheimnisvolle Ruinen, kämpft gegen gefährliche Gegner, helft Bewohnern eines Dorfes, entdeckt verborgene Schätze und erzählt gemeinsam eure eigene Geschichte.

Anders als bei vielen anderen Spielen gibt es dabei keinen Gewinner und keinen Verlierer. Statt gegeneinander zu spielen, arbeitet ihr als Gruppe zusammen und versucht gemeinsam, die Herausforderungen eures Abenteuers zu meistern.

Dieses Kapitel erklärt die Grundlagen von Heldenpfade und zeigt euch, wie ihr sofort mit eurem ersten Abenteuer beginnen könnt.

\section{Was ist Heldenpfade?}

Heldenpfade ist ein Pen-\&-Paper-Rollenspiel.

Gemeinsam erzählt ihr eine Geschichte und erlebt darin ein Abenteuer. Ein Spieler übernimmt die Spielleitung und erzählt das Abenteuer. Alle anderen Spieler übernehmen jeweils einen Helden und entscheiden selbst, was ihre Helden denken, sagen und tun.

Die Spieler sprechen miteinander, planen gemeinsam und versuchen als Gruppe, die Herausforderungen ihres Abenteuers zu meistern.

Die Spielleitung beschreibt die Welt, übernimmt alle anderen Figuren und entscheidet, wie die Umwelt auf die Handlungen der Helden reagiert. Sie spielt dabei nicht gegen die Spieler. Ihre Aufgabe ist es, gemeinsam mit allen am Tisch eine spannende Geschichte zu erzählen und dafür zu sorgen, dass alle Spaß am Spiel haben.

Dabei gibt es fast keine festen Grenzen. Vielleicht möchtet ihr mit einem Händler verhandeln, einen steilen Felsen erklimmen, ein Rätsel lösen oder gegen einen Drachen kämpfen. Fast jede Idee ist erlaubt.

Ob eine Handlung gelingt, entscheiden meistens die Regeln und die Würfel. Sie sorgen dafür, dass Erfolg und Misserfolg spannend bleiben und niemand vorher weiß, wie die Geschichte ausgeht.

Die Regeln sollen eure Fantasie dabei unterstützen -- sie sollen sie nicht einschränken.

\section{Wie läuft ein Spiel ab?}

Eine Partie Heldenpfade besteht aus einem ständigen Wechsel zwischen Erzählen, Entscheiden und Würfeln.

Zunächst beschreibt die Spielleitung eine Situation. Danach entscheiden die Spieler, wie ihre Helden darauf reagieren.

Ist unklar, ob eine Handlung gelingt, werden Würfel verwendet. Das Ergebnis entscheidet, wie die Geschichte weitergeht.

Ein kurzer Ausschnitt könnte zum Beispiel so aussehen:

\begin{quote}
\textbf{Spielleitung:} „Vor euch steht eine alte Holztür. Dahinter hört ihr leise Stimmen.“

\textbf{Spieler:} „Ich lausche zuerst an der Tür.“

\textbf{Spielleitung:} „Dann führe bitte eine Wahrnehmungsprobe durch.“
\end{quote}

Der Spieler würfelt und erfährt anschließend, was sein Held hört.

So entwickelt sich die Geschichte Schritt für Schritt weiter.

\section{Kann ich sofort losspielen?}

Ja.

Heldenpfade enthält mehrere spielfertige Beispielhelden. Ihr könnt euch jeweils einen dieser Helden aussuchen und sofort mit eurem ersten Abenteuer beginnen.

Die Beispielhelden befinden sich in Anhang 5 -- Beispielhelden.

Wenn ihr lieber einen eigenen Helden erschaffen möchtet, findet ihr die vollständigen Regeln dazu in Kapitel 9 -- Wie entsteht ein Held?

\section{Was brauche ich zum Spielen?}

Zum Spielen von Heldenpfade benötigt ihr:

\begin{itemize}
  \item dieses Regelwerk,
  \item ein Abenteuer,
  \item Charakterbögen mit Beispielhelden,
  \item Karten,
  \item Marker,
  \item Stifte
  \item sowie die zum Spiel gehörenden Würfel.
\end{itemize}

\section{Welche Würfel gibt es?}

In Heldenpfade kommen verschiedene Würfelarten zum Einsatz. Die Symbole werden im gesamten Regelwerk immer gleich verwendet.

\begin{center}
\begin{hptable}{C{1.7cm} L{4.0cm} Y}
\textbf{Symbol} & \textbf{Würfelart} & \textbf{Verwendung} \\
\heroDie{} & Heldenwürfel & Der persönliche Würfel eines Helden. Jeder Held besitzt genau einen Heldenwürfel. \\
\greenDie{} & Attributswürfel & Grüne Würfel für die natürlichen Eigenschaften und grundlegenden Stärken eines Helden. \\
\blueDie{} & Fertigkeitswürfel & Blaue Würfel für erlernte, geübte und trainierte Fähigkeiten. \\
\orangeDmgDie{} & Leichter Schadenswürfel & Orangener Würfel für leichten Schaden. \\
\redDmgDie{} & Mittlerer Schadenswürfel & Roter Würfel für mittleren Schaden. \\
\blackDmgDie{} & Schwerer Schadenswürfel & Schwarzer Würfel für schweren Schaden. \\
\end{hptable}
\end{center}

Welche Würfel bei einer Probe verwendet werden, hängt von den Fähigkeiten des Helden ab. Je gefährlicher ein Angriff ist, desto stärkere Schadenswürfel können zum Einsatz kommen.

Die genaue Verwendung aller Würfel wird in den folgenden Kapiteln erklärt.

\section{Wie ist dieses Regelwerk aufgebaut?}

Das Regelwerk ist so aufgebaut, dass ihr die Spielregeln Schritt für Schritt kennenlernt.

Die ersten Kapitel erklären die wichtigsten Grundlagen des Spiels. Danach folgen die Regeln für Kampf, Ausrüstung, Zauber und Talente. Anschließend erfahrt ihr, wie ihr euren eigenen Helden erschafft und weiterentwickelt.

Wenn ihr später eine bestimmte Regel nachschlagen möchtet, könnt ihr direkt zum passenden Kapitel springen.

\chapter{Was macht einen Helden aus?}

Ein Held wird durch wenige zentrale Bausteine beschrieben: Attribute, Fertigkeiten, Lebenspunkte, Mana, Ausrüstung und Talente.

Diese Werte zeigen nicht nur, was ein Held kann. Sie zeigen vor allem, wie er Herausforderungen löst: mutig, klug, geschickt, zäh, überzeugend, flink oder mit besonderen Fähigkeiten.

Dieses Kapitel gibt einen Überblick über diese Bausteine. Die genauen Regeln für Proben, Kämpfe, Ausrüstung, Mag
\chapter{Wann gelingt ein Vorhaben?}
\label{chap:checks}

Helden wollen ständig etwas tun: über einen Bach springen, ein Schloss öffnen, einen Streit schlichten, eine Spur finden oder sich an einem schlafenden Troll vorbeischleichen.

Manchmal ist sofort klar, was passiert. Dann erzählt ihr einfach weiter. Manchmal ist aber unklar, ob ein Vorhaben gelingt. Dann wird gewürfelt. Das nennt man eine Probe.

\section{Wann wird gewürfelt?}

\begin{kurzregel}[title=Wann ist eine Probe nötig?]
Eine Probe wird nur durchgeführt, wenn
\begin{itemize}
  \item das Vorhaben grundsätzlich möglich ist,
  \item der Ausgang unklar ist und
  \item Erfolg oder Fehlschlag einen Unterschied machen.
\end{itemize}
\end{kurzregel}

Ist ein Vorhaben leicht und spricht nichts dagegen, gelingt es ohne Probe. Ist es unmöglich, wird ebenfalls nicht gewürfelt.

\section{Was ist eine Probe?}

Eine Probe ist ein Würfelwurf, mit dem geprüft wird, ob ein Vorhaben gelingt.

Bei Proben werden vier grundlegende Symbole verwendet:

\begin{itemize}
  \item der \textbf{Heldenwürfel (\heroDie)},
  \item grüne \textbf{Attributswürfel (\greenDie)},
  \item blaue \textbf{Fertigkeitswürfel (\blueDie)},
  \item der \textbf{Stern (\star)} als Erfolgssymbol.
\end{itemize}

\begin{kurzregel}[title=Grundregel einer Probe]
\textbf{Würfelpool:} 1\heroDie{} + passende \greenDie{} + passende \blueDie

\textbf{Erfolg:} erzielte \star{} $\geq$ Schwierigkeit
\end{kurzregel}

Jede Probe eines Helden enthält genau einen \heroDie. Dazu kommen die \greenDie{} der passenden Attribute und die \blueDie{} der passenden Fertigkeit.

Nach dem Wurf werden alle \star{} zusammengezählt. Erreicht der Held mindestens die von der Spielleitung festgelegte Schwierigkeit, gelingt die Probe.

\section{Den Würfelpool bilden}

Jede Fertigkeit gehört zu zwei Attributen. Für eine Fertigkeitsprobe nimmt der Spieler:

\begin{itemize}
  \item \heroDie,
  \item die \greenDie{} beider zugehöriger Attribute,
  \item die \blueDie{} der verwendeten Fertigkeit.
\end{itemize}

\begin{beispiel}[title=Würfelpool einer Fertigkeitsprobe]
Turnen gehört zu GEW und KRA. Kiera besitzt GEW 2, KRA 1 und Turnen 1.

Sie würfelt daher:

\textbf{1\heroDie{} + 3\greenDie{} + 1\blueDie}
\end{beispiel}

\subsection{Attributsproben}

Passt keine Fertigkeit, kann die Spielleitung eine Attributsprobe verlangen.

\begin{kurzregel}[title=Attributsprobe]
Für eine Attributsprobe verwendet der Held
\begin{itemize}
  \item immer 1\heroDie{} und
  \item entweder die \greenDie{} zweier passender Attribute oder die \greenDie{} eines besonders wichtigen Attributs zweimal.
\end{itemize}
\end{kurzregel}

\begin{beispiel}[title=Ein Attribut doppelt verwenden]
Kiera versucht, eine schwere Tür aufzudrücken. Die Spielleitung entscheidet, dass allein ihre Kraft entscheidend ist.

Kiera besitzt KRA 3. Da Kraft doppelt verwendet wird, würfelt sie:

\textbf{1\heroDie{} + 3\greenDie{} + 3\greenDie}
\end{beispiel}

\section{Eine Probe in sieben Schritten}

\begin{enumerate}
  \item Der Spieler beschreibt das Vorhaben.
  \item Die Spielleitung entscheidet, ob eine Probe nötig ist.
  \item Die Spielleitung bestimmt Fertigkeit oder Attribute und nennt die Schwierigkeit.
  \item Der Spieler bildet den Würfelpool und würfelt.
  \item Vorteil, Nachteil und Mondkristalle werden angewendet.
  \item Alle \star{} werden gezählt und mit der Schwierigkeit verglichen.
  \item Die Spielleitung beschreibt Erfolg, Fehlschlag oder Teilerfolg.
\end{enumerate}

\section{Schwierigkeit}

\begin{hptable}{L{3cm} C{2.5cm} Y}
\textbf{Schwierigkeit} & \textbf{Benötigte \star} & \textbf{Beispiel} \\
Leicht & 1 & Über einen breiten Bach springen \\
Normal & 2 & Sich an einer Wache vorbeischleichen \\
Schwierig & 3 & Einen verschlossenen Tresor knacken \\
Sehr schwierig & 4 & Eine außergewöhnlich anspruchsvolle Aufgabe unter Zeitdruck bewältigen \\
Heldentat & 5 & Etwas vollbringen, das selbst für erfahrene Helden beinahe unmöglich ist \\
\end{hptable}

Die Spielleitung legt die Schwierigkeit passend zur Situation fest. Die Beispiele dienen nur als Orientierung; Umstände, Zeitdruck, Vorbereitung und verfügbare Hilfsmittel können dieselbe Aufgabe leichter oder schwerer machen.

\section{Erfolg, Fehlschlag und Teilerfolg}

Erreicht der Held genügend \star, gelingt das Vorhaben. Bei zu wenigen \star{} misslingt es oder gelingt nur teilweise.

Ein Fehlschlag soll die Geschichte nicht anhalten. Er verändert die Situation: Der Held verliert Zeit, macht Lärm, gerät in Gefahr, verliert etwas oder braucht einen anderen Plan.

Ein Teilerfolg bedeutet, dass das Vorhaben gelingt, aber ein Problem entsteht.

\begin{beispiel}[title=Teilerfolg]
Kiera knackt ein Schloss, erreicht aber nicht genügend \star{} für einen vollständigen Erfolg. Die Tür öffnet sich trotzdem, doch das Werkzeug bricht ab und verursacht so viel Lärm, dass eine Wache aufmerksam wird.
\end{beispiel}

\section{Vorteil und Nachteil}

\begin{kurzregel}[title=Vorteil und Nachteil]
Vorteil und Nachteil gelten nur für Proben von Helden, nicht für Gegner oder Schadenswürfe.

Treffen Vorteil und Nachteil gleichzeitig zu, heben sie sich gegenseitig auf.
\end{kurzregel}

\subsection{Vorteil}

Bei Vorteil darf der Spieler beliebige Würfel dieser Probe einmal kostenlos neu würfeln. Danach gilt das neue Ergebnis.

\subsection{Nachteil}

Bei Nachteil darf der Spieler für diese Probe keine Mondkristalle einsetzen.

\section{Monde und Mondkristalle}

\begin{kurzregel}[title=Monde]
Ein \moon{} zählt nicht als \star.

Zeigt eine Würfelseite sowohl \moon{} als auch \star, wird der \star{} trotzdem normal gezählt.
\end{kurzregel}

Jeder gewürfelte \moon{} erzeugt sofort einen \textbf{Mondkristall (\crystal)}. Effekte oder besondere Würfel können später auch mehrere Mondkristalle erzeugen. Der Spieler kann erhaltene \crystal{} noch in derselben Probe oder später einsetzen. Ein Held kann beliebig viele \crystal{} besitzen.

Wann ungenutzte \crystal{} verfallen, bestimmt die Spielleitung passend zum Abenteuer.

\begin{todo}[title={TODO \#1 -- Regeldesign: Mondkristalle außerhalb einer Probe}]
\textbf{Offene Designentscheidung}

Der Umgang mit Mondkristallen außerhalb einer Probe ist noch nicht endgültig festgelegt.

Zu klären:
\begin{itemize}
  \item Gibt es ein Maximum für gespeicherte Mondkristalle?
  \item Wann verfallen ungenutzte Mondkristalle?
  \item Bleiben Mondkristalle zwischen Szenen, Begegnungen oder Abenteuern erhalten?
\end{itemize}
\end{todo}

\subsection{Mondkristalle einsetzen}

\begin{kurzregel}[title=Mondkristall einsetzen]
Gib 1\crystal{} aus und würfle alle Würfel \textbf{einer Farbe} neu:
\begin{itemize}
  \item den \heroDie,
  \item alle \greenDie{} oder
  \item alle \blueDie.
\end{itemize}
Danach gilt das neue Ergebnis.
\end{kurzregel}

Mehrere \crystal{} können in derselben Probe eingesetzt werden. Würfel, die bereits für einen anderen Effekt verwendet wurden, sind fest und dürfen in dieser Probe nicht mehr neu gewürfelt werden.

Talente und andere Effekte können zusätzliche Möglichkeiten bieten, \crystal{} auszugeben.

\section{Ausführliches Beispiel}

\begin{beispiel}[title=Eine schwierige Turnen-Probe]
\textbf{1. Würfelpool}

Kiera besitzt GEW 2, KRA 1 und Turnen 1. Sie würfelt daher \heroDie{} + 3\greenDie{} + 1\blueDie.

\textbf{2. Schwierigkeit}

Die Spielleitung legt die Schwierigkeit auf \textbf{schwierig} fest. Kiera benötigt also 3\star.

\textbf{3. Erster Wurf}

Der \heroDie{} zeigt einen \moon. Zwei andere Würfel zeigen zusammen 2\star.

\textbf{4. Mondkristall}

Der gewürfelte \moon{} erzeugt sofort 1\crystal.

\textbf{5. Neuwurf}

Kiera gibt den \crystal{} aus und würfelt alle \greenDie{} neu. Der \heroDie{} und der \blueDie{} bleiben liegen.

\textbf{6. Ergebnis}

Nach dem Neuwurf erzielt Kiera insgesamt 3\star{} und besteht die Probe.
\end{beispiel}
\chapter{Wie funktioniert ein Kampf?}
\label{chap:combat}

Ein Kampf ist eine besondere Szene, in der die Zeit in Runden und Züge aufgeteilt wird.

\section{Die Kampfrunde}

In jeder Runde handeln zuerst alle Helden, danach alle Gegner. Die Helden bestimmen ihre Reihenfolge selbst. Sobald alle Gegner vollständig aktiviert wurden, endet die Runde und die nächste beginnt.

\begin{kurzregel}[title=Reihenfolge einer Kampfrunde]
\textbf{Helden handeln} $\rightarrow$ \textbf{Gegner handeln}
\end{kurzregel}

\subsection{Der Zug eines Helden}

Ein Held besitzt in seinem Zug normalerweise:

\begin{itemize}
  \item eine \textbf{Standardaktion (\standardAction)},
  \item eine \textbf{Bewegungsaktion (\moveAction)},
  \item beliebig viele \textbf{freie Aktionen (\freeAction)}.
\end{itemize}

Mit einer \standardAction{} kann ein Held beispielsweise angreifen, einen Zauber wirken oder eine besondere Fähigkeit einsetzen.

Mit einer \moveAction{} kann er sich bewegen oder seine Ausrüstung neu anordnen. Seine \standardAction{} darf er stattdessen als zweite \moveAction{} verwenden.

\freeAction{} sind kurze Handlungen wie etwas rufen, auf etwas zeigen oder einen Gegenstand fallen lassen. Sie dürfen keine umfangreiche Handlung ersetzen.

Wird eine Aktion beschleunigt, wird eine \standardAction{} zur \moveAction{} und eine \moveAction{} zur \freeAction. Eine \freeAction{} kann nicht weiter beschleunigt werden.

\begin{tip}[title=Tipp: Aktionen zuerst planen]
Kläre zu Beginn deines Zuges kurz, was du mit deiner Standardaktion und deiner Bewegungsaktion erreichen möchtest. Das beschleunigt den Zug und verhindert unnötige Rückfragen während der Ausführung.
\end{tip}

\section{Bewegung, Nähe und Entfernung}

Figuren können zueinander \textbf{benachbart} oder \textbf{entfernt} sein. Eine Figur kann gleichzeitig zu mehreren Figuren benachbart sein.

Wie viele Bewegungen nötig oder mit einer \moveAction{} möglich sind, ergibt sich aus der Battlemap oder aus der Beschreibung der Spielleitung. Hindernisse und die konkrete Situation können eine Bewegung erschweren oder mehrere Bewegungen erforderlich machen.

Ein Held darf sich von einem benachbarten Gegner entfernen, ohne dass dies automatisch eine weitere Folge auslöst. Gegnerfähigkeiten können davon abweichen.

\begin{todo}
\textbf{Regeldesign: Bewegung im Kampf}

Die abstrakte Bewegung über \emph{benachbart} und \emph{entfernt} muss später mit den Battlemap-Regeln abgestimmt werden. Zu klären ist insbesondere, wie viele Felder oder Zonen eine Bewegungsaktion umfasst und wie Hindernisse behandelt werden.
\end{todo}

\section{Ausrüstung im Kampf wechseln}

Mit einer \moveAction{} darf ein Held seine ausgerüsteten und verstauten Gegenstände beliebig neu anordnen, solange danach alle Ausrüstungsplätze korrekt belegt sind.

Der Gegenstand am Körper-Ausrüstungsplatz ist davon ausgenommen. Eine Körperrüstung anzulegen, abzulegen oder zu wechseln, benötigt eine eigene \standardAction. Die allgemeinen Regeln für Ausrüstungsplätze und das Wechseln von Ausrüstung stehen in \chapterref{chap:equipment}.

\section{Einen Gegner angreifen}

Ein Angriff benötigt eine \standardAction. \textbf{Angriffsoptionen} beschreiben, wie eine Figur einen Angriff ausführt. Sie befinden sich meist auf \textbf{Waffen}, können aber auch auf \textbf{Zaubern} stehen. Vereinzelt verleihen auch andere \textbf{Ausrüstungsgegenstände} oder \textbf{Talente} eigene Angriffsoptionen. Die zugehörigen Regeln stehen in \chapterref{chap:equipment}, \chapterref{chap:magic} und \chapterref{chap:talents}.

Jede Angriffsoption legt fest,

\begin{itemize}
  \item welche Angriffsfertigkeit verwendet wird,
  \item welche Kaufoptionen zur Verfügung stehen und
  \item welche zusätzlichen Folgen nach dem Angriff eintreten.
\end{itemize}

Eine Angriffsoption auf einer Ausrüstungskarte kann nur gewählt werden, wenn sie auf der aktuell sichtbaren Kartenseite abgebildet ist.

\begin{kurzregel}[title=Einen Angriff durchführen]
\begin{enumerate}
  \item Angriffsoption wählen.
  \item Zulässiges Ziel bestimmen.
  \item Angriffsprobe ablegen.
  \item Bei mindestens einem \star{} Kaufoptionen erwerben.
  \item Verbleibende \star{} zum Verbessern erhaltener Schadenswürfel ausgeben.
  \item Schadenspools durch Effekte verändern und gegebenenfalls vor dem Wurf auf mehrere Ziele verteilen.
  \item Jeden Schadenspool getrennt werfen und den Schaden addieren.
  \item Rüstungsbruch und \armor{} für jedes Ziel anwenden.
  \item Verbleibenden Schaden von den jeweiligen Lebenspunkten abziehen.
  \item Gekaufte Kampfeffekte und zusätzliche Folgen zu ihrem jeweils angegebenen Zeitpunkt abhandeln.
\end{enumerate}
\end{kurzregel}

Für die Angriffsfertigkeiten gelten folgende Reichweiten:

\begin{itemize}
  \item \textbf{Schlagen} kann benachbarte Gegner als Ziel haben.
  \item \textbf{Schießen} kann entfernte Gegner als Ziel haben. Ein Held darf auch schießen, während ein Gegner zu ihm benachbart ist; das Ziel des Angriffs muss dennoch entfernt sein.
  \item \textbf{Zaubern} kann benachbarte oder entfernte Gegner als Ziel haben, sofern der Zauber nichts anderes bestimmt.
\end{itemize}

\subsection{Treffen}

Der Spieler legt mit der angegebenen Angriffsfertigkeit eine Angriffsprobe nach den Regeln aus \chapterref{chap:checks} ab. Zeigt sie keinen \star, verfehlt der Angriff. Jeder erzielte \star{} kann anschließend ausgegeben werden, um eine oder mehrere Kaufoptionen der gewählten Angriffsoption zu erwerben. Die vollständigen Regeln für Kaufoptionen stehen in \chapterref{chap:equipment}.

Helden erhalten Vorteil auf Angriffsproben gegen Gegner mit dem Statuseffekt \emph{Zu Boden}. Ein Held mit dem Statuseffekt \emph{Zu Boden} hat Nachteil auf seine eigenen Angriffsproben.

Gegner weichen Angriffen von Helden grundsätzlich nicht aus.

\subsection{Kaufoptionen und Schaden}

Es gibt drei Arten von Schadenswürfeln:

\begin{itemize}
  \item den \textbf{orangen Schadenswürfel (\orangeDmgDie)} für leichten Schaden,
  \item den \textbf{roten Schadenswürfel (\redDmgDie)} für mittleren Schaden,
  \item den \textbf{schwarzen Schadenswürfel (\blackDmgDie)} für schweren Schaden.
\end{itemize}

Eine Kaufoption nennt ihre Kosten in \star{} und den dafür erhaltenen Effekt. Jede Kaufoption der gewählten Angriffsoption darf pro Angriff höchstens einmal gekauft werden. Der Angreifer darf mehrere unterschiedliche Kaufoptionen in beliebiger Kombination kaufen, solange er ihre Kosten bezahlt und alle Voraussetzungen erfüllt.

Nachdem alle gewünschten Kaufoptionen gekauft wurden, dürfen verbleibende \star{} bereits erhaltene Schadenswürfel verbessern. Jeder ausgegebene \star{} verbessert einen Würfel um eine Stufe: \orangeDmgDie{} wird zu \redDmgDie{}, \redDmgDie{} wird zu \blackDmgDie{}. Derselbe Würfel darf mehrfach verbessert werden; \blackDmgDie{} können nicht weiter verbessert werden. Danach verfallen noch übrige \star{}.

Alle gekauften und verbesserten Schadenswürfel bilden einen gemeinsamen Schadenspool. Sie werden gemeinsam geworfen und ihre Ergebnisse anschließend addiert. Kaufoptionen können statt Schadenswürfeln oder zusätzlich zu ihnen Kampfeffekte verleihen.

\begin{kurzregel}[title=Schaden bestimmen]
\begin{enumerate}
  \item Kaufoptionen jeweils höchstens einmal kaufen.
  \item Alle erhaltenen Schadenswürfel zu einem Schadenspool zusammenfügen.
  \item Verbleibende \star{} ausgeben, um Schadenswürfel stufenweise zu verbessern.
  \item Den Schadenspool durch weitere Effekte verändern.
  \item Falls ein Effekt wie Spalten dies erlaubt, die Würfel vor dem Wurf auf Ziele verteilen.
  \item Jeden Schadenspool getrennt werfen und den gewürfelten Schaden addieren.
  \item Rüstungsbruch und Schutz durch \armor{} für jedes Ziel anwenden.
  \item Verbleibenden Schaden als Verlust von Lebenspunkten anwenden.
  \item Weitere Kampfeffekte zu ihrem angegebenen Zeitpunkt abhandeln.
\end{enumerate}
\end{kurzregel}

\begin{beispiel}[title=Beispiel: Schadenswürfel kaufen und verbessern]
Ein Schwert besitzt die Kaufoptionen 2\star{} $\rightarrow$ \orangeDmgDie{}\orangeDmgDie{} und 3\star{} $\rightarrow$ \orangeDmgDie{}\orangeDmgDie{}\orangeDmgDie{}. Mit 6\star{} kauft ein Held beide Optionen je einmal. Für den verbleibenden \star{} verbessert er einen der fünf orangen Würfel zu einem roten Würfel. Sein Schadenspool besteht aus 1\redDmgDie{} und 4\orangeDmgDie{}.
\end{beispiel}

\begin{beispiel}[title=Beispiel: Schaden und Kampfeffekt kombinieren]
Eine Axt besitzt unter anderem die Kaufoptionen 2\star{} $\rightarrow$ \redDmgDie{} und 2\star{} $\rightarrow$ \effectArmorBreak{} 2. Mit 4\star{} kann ein Held beide Optionen kaufen. Der Angriff verursacht Schaden mit 1\redDmgDie{} und erhält zusätzlich 2 Rüstungsbruch. Beide Kaufoptionen sind damit für diesen Angriff verbraucht.
\end{beispiel}

Das \textbf{Schild-Symbol (\armor)} beschreibt Schutz durch Rüstung, Schilde, Talente oder andere Effekte. Jedes \armor{} verhindert 1 Schaden. \textbf{Rüstungsbruch (\armorPierce)} ignoriert die angegebene Anzahl an Schild-Symbolen des Ziels, unabhängig von ihrer Quelle.

\begin{kurzregel}[title=Schild-Symbole und Rüstungsbruch]
Jedes \armor{} verhindert \textbf{1 Schaden}. Jedes \armorPierce{} ignoriert \textbf{1 \armor{}} des Ziels.
\end{kurzregel}

Der verbleibende Schaden lässt die Figur entsprechend viele \textbf{Lebenspunkte verlieren (\lifeLoss)}.

\section{Wenn Gegner angreifen}

Gegner legen keine Angriffsprobe ab. Ihr Profil gibt unmittelbar die verwendeten Schadenswürfel an.

Bevor diese Würfel geworfen werden, darf der angegriffene Held gegebenenfalls \textbf{ausweichen (\evade)}. Welche Würfel er für seine Ausweichprobe verwendet, ergibt sich aus seiner Ausrüstung oder seinen Talenten.

\begin{kurzregel}[title=Ausweichen]
Ein Held kann nur ausweichen, wenn seine \textbf{Ausrüstung oder ein Talent} dies erlaubt.

Jeder \star{} der Ausweichprobe entfernt einen beliebigen Schadenswürfel des Gegners. Überschüssige \star{} verfallen. Entfernte Schadenswürfel werden nicht geworfen.
\end{kurzregel}

Ein Held mit dem Statuseffekt \emph{Zu Boden} hat Nachteil auf seine Ausweichprobe. Greift ein Gegner mit dem Statuseffekt \emph{Zu Boden} an, hat der angegriffene Held Vorteil auf seine Ausweichprobe. Kann der Held aufgrund seiner Ausrüstung oder Talente nicht ausweichen, entfällt die Probe vollständig.

Die verbleibenden Schadenswürfel werden geworfen. Anschließend verhindert jedes \armor{} des Helden 1 Schaden. Für den übrigen Schaden verliert der Held entsprechend viele Lebenspunkte \lifeLoss.

\begin{todo}
\textbf{Balancing: Wahl der Schadenswürfel beim Ausweichen}

Aktuell darf der Held bei einem gemischten Schadenspool frei wählen, welche Schadenswürfel durch seine Ausweicherfolge entfernt werden. Prüfen, ob diese freie Wahl dauerhaft sinnvoll und ausgewogen ist.
\end{todo}

\section{Kampfunfähig oder besiegt}

Lebenspunkte können niemals unter 0 sinken.

Sinkt ein Gegner auf 0 Lebenspunkte, ist er besiegt.

Sinkt ein Held auf 0 Lebenspunkte, ist er kampfunfähig und kann nicht handeln. Wird er geheilt und erhält mindestens 1 \textbf{Lebenspunkt (\lifeGain)}, ist er wieder kampfbereit, erhält aber den Statuseffekt \emph{Zu Boden}. Er muss daher zunächst entsprechend den Regeln in \appendixref{app:status-effects} aufstehen.

Wenn alle Helden kampfunfähig sind, sind die Helden besiegt. Das bedeutet nicht automatisch, dass sie sterben.

\section{Wann endet ein Kampf?}

Ein Kampf endet, wenn die gefährliche Auseinandersetzung vorbei ist, beispielsweise weil alle Gegner besiegt sind, eine Seite flieht oder aufgibt oder das Ziel des Kampfes erreicht wurde.

Wenn der Kampf endet, wird wieder normal weitererzählt.

\section{Beispiele}

\begin{beispiel}[title=Beispiel: Ein Held greift an]
\begin{enumerate}
  \item \textbf{Angriffsoption wählen (Zahnschwert):} Kiera greift einen benachbarten Gegner mit ihrem Zahnschwert an.
  \item \textbf{Angriffsprobe:} Sie legt die angegebene Schlagen-Probe ab und erzielt 6\star.
  \item \textbf{Kaufoptionen erwerben:} Kiera kauft für 2\star{} zwei \redDmgDie{} und für 3\star{} drei \orangeDmgDie{}.
  \item \textbf{Schadenswürfel verbessern:} Mit dem verbleibenden \star{} verbessert sie einen \orangeDmgDie{} zu einem \redDmgDie{}.
  \item \textbf{Schaden würfeln:} Ihr Schadenspool aus 3\redDmgDie{} und 2\orangeDmgDie{} verursacht zusammen 8 Schaden.
  \item \textbf{Schutz anwenden:} Der Gegner besitzt 2\armor{} und verhindert damit 2 Schaden.
  \item \textbf{Lebenspunkte verlieren:} Der Gegner verliert die verbleibenden 6 Lebenspunkte.
\end{enumerate}
\end{beispiel}

\begin{beispiel}[title=Beispiel: Ein Gegner greift an]
\begin{enumerate}
  \item \textbf{Schadenspool bestimmen:} Ein Gegner greift Kiera mit 1\blackDmgDie{} und 2\orangeDmgDie{} an.
  \item \textbf{Ausweichen:} Kieras Ausrüstung erlaubt eine Ausweichprobe. Sie erzielt 1\star{} und entfernt den \blackDmgDie.
  \item \textbf{Schaden würfeln:} Die beiden verbleibenden \orangeDmgDie{} verursachen zusammen 3 Schaden.
  \item \textbf{Schutz anwenden:} Kiera besitzt 1\armor{} und verhindert damit 1 Schaden.
  \item \textbf{Lebenspunkte verlieren:} Kiera verliert die verbleibenden 2 Lebenspunkte.
\end{enumerate}
\end{beispiel}

\chapter{Wie funktioniert Ausrüstung?}

\section{Was ist Ausrüstung?}

Ausrüstung sind Gegenstände, die ein Held bei sich trägt und verwenden kann. Sie wird durch Karten dargestellt.

\section{Ausrüstungsplätze}

Ein Held besitzt die Ausrüstungsplätze Kopf, zwei Hände, Körper, Gürtel und Füße. Nur Gegenstände, die an ihrem passenden Platz liegen, gelten als ausgerüstet und können verwendet werden.

Ein Held kann außerdem beliebig viel verstaute Ausrüstung mit sich tragen. Verstaute Ausrüstung gilt nicht als ausgerüstet und ihre Angaben und Effekte können nicht genutzt werden.

\subsection{Hände}

Ein Held hat zwei Hände. Manche Gegenstände benötigen eine Hand, andere beide Hände. Die ausgerüsteten Gegenstände dürfen zusammen höchstens zwei Hände belegen.

\section{Ausrüstung wechseln}

Mit einer \moveAction{} darf ein Held seine ausgerüsteten und verstauten Gegenstände beliebig neu anordnen. Er kann dabei mehrere Ausrüstungsplätze gleichzeitig verändern.

Der Körperplatz ist eine Ausnahme: Einen Gegenstand am Körperplatz auszurüsten, abzulegen oder gegen einen anderen zu tauschen kostet eine eigene \standardAction.

\section{Ausrüstungskarten lesen}

Waffen besitzen eine oder mehrere Angriffsoptionen. Jede Angriffsoption nennt die Angriffsfertigkeit, die Umrechnung von \star{} in Schadenswürfel und mögliche besondere Effekte.

Rüstungen und andere Gegenstände können eine Ausweichprobe (\evade) erlauben oder zusätzliche \blueDie{} dafür gewähren. Trägt ein Held keine Rüstung am Körperplatz, darf er mit einer Turnen-Probe ausweichen.

Schutz wird durch Schild-Symbole (\armor) dargestellt. Jedes \armor{} verhindert 1 Schaden. Rüstungsdurchdringung (\armorPierce) ignoriert die angegebene Anzahl an \armor{} des Ziels.

\section{Erschöpfte Karten}

Manche Effekte fordern den Spieler auf, eine Karte zu \textbf{erschöpfen (\exhaust)}. Dazu wird die Karte umgedreht.

Eine erschöpfte Karte bleibt ausgerüstet und belegt weiterhin ihren Ausrüstungsplatz. Da ihre Vorderseite nicht sichtbar ist, können jedoch keine ihrer Angaben, Werte oder Effekte genutzt werden.

Erschöpfte Karten werden spätestens bei der nächsten Rast wieder bereitgemacht. Andere Regeln können sie früher bereitmachen.

\section{Weiterführende Informationen}

Im Anhang befinden sich die vollständige Ausrüstungsliste und die Symbolübersicht.

\chapter{Wie wirkt Magie?}
\label{chap:magic}

\section{Was ist Magie?}

Magie gehört zu den außergewöhnlichsten Fähigkeiten, die ein Held in Heldenpfade erlernen kann. Magie kann Verletzungen heilen, Feuer beschwören, Verbündete stärken oder Gegner schwächen.

Nicht jeder Held kann zaubern. Ein Held benötigt das Talent \emph{Magiebegabt}, einen erlernten und verfügbaren Zauber sowie ausreichend Mana, um dessen Kosten zu bezahlen.

\section{Mana}

Jeder Zauber kostet Mana. Beim ersten Auftreten wird das Ausgeben von Mana als \textbf{Mana bezahlen (\manaLoss)} und das Erhalten von Mana als \textbf{Mana erhalten (\manaGain)} bezeichnet.

Jeder Held besitzt einen Manavorrat und aktuelles Mana. Der Manavorrat ist das Maximum, das aktuelle Mana steht für die momentan verfügbare Menge.

Zu Beginn jedes Abenteuers besitzt ein Held seinen vollständigen Manavorrat. Reicht das aktuelle Mana nicht aus, kann ein Zauber nicht gewirkt werden. Das aktuelle Mana kann niemals unter 0 oder über den Manavorrat steigen.

Verbrauchtes Mana wird während einer Rast mit \manaGain{} regeneriert.

\section{Magiebegabt}

Um Zauber wirken zu können, muss ein Held das Talent \emph{Magiebegabt} besitzen. Das Talent erhöht den Manavorrat des Helden um 6.

\section{Zauber und ihre Effekte}

Jeder Zauber wird durch eine Zauberkarte beschrieben. Die Karte nennt den Namen, die Manakosten, die Zauberzeit und den Effekt des Zaubers. Sie kann außerdem Ziele, eine Wirkungsdauer, Schlüsselwörter oder weitere Voraussetzungen angeben.

Ein \textbf{Zauber} ist die gesamte Karte. Um ihn zu wirken, bezahlt der Held seine Manakosten und führt den beschriebenen \textbf{Effekt} aus. Ein Effekt kann eine \textbf{Angriffsoption} sein. In diesem Fall bestimmt die Angriffsoption wie gewohnt die Angriffsfertigkeit, wie erzielte \star{} ausgegeben werden können, und mögliche zusätzliche Folgen. Jede aufgeführte Zeile darf pro Angriff höchstens einmal genutzt werden, sofern sie nicht ausdrücklich als mehrfach nutzbar bezeichnet ist; verbleibende \star{} können anschließend erhaltene Schadenswürfel verbessern.

Zauberkarten werden durch das Wirken eines Zaubers nicht automatisch erschöpft. Nur wenn eine Zauberkarte ausdrücklich das Symbol \textbf{Erschöpfen (\exhaust)} zeigt oder der Regeltext dies verlangt, wird genau diese Karte umgedreht. Danach gelten ausschließlich die Angaben und Effekte der nun sichtbaren Kartenseite.

Ein Held kann einen erlernten Zauber wiederholt wirken, solange genügend Mana vorhanden ist, die aktuell sichtbare Kartenseite das Wirken erlaubt und alle Voraussetzungen erfüllt sind.

\section{Einen Zauber wirken}

Jeder Zauber nennt eine Zauberzeit. Im Kampf gibt sie an, welche Aktion zum Wirken benötigt wird. Ein besonders schneller Zauber kann beispielsweise als \freeAction{} gewirkt werden, während ein aufwendigerer Zauber eine \standardAction{} benötigt.

\begin{kurzregel}[title=Einen Zauber wirken]
\begin{enumerate}
  \item Zauber wählen.
  \item Angegebene Manakosten mit \manaLoss{} bezahlen.
  \item Gültige Ziele bestimmen.
  \item Effekt vollständig ausführen.
  \item Falls angegeben, die Zauberkarte mit \exhaust{} erschöpfen.
\end{enumerate}
\end{kurzregel}

Kann ein Schritt nicht vollständig ausgeführt werden, kann der Zauber nicht gewirkt werden. Insbesondere müssen ausreichend Mana und mindestens ein gültiges Ziel vorhanden sein.

\begin{beispiel}[title=Beispiel: Einen unterstützenden Zauber wirken]
Lina möchte mit dem Zauber \emph{Heilender Funke} einen benachbarten Verbündeten heilen. Der Zauber kostet 1\manaLoss{} und benötigt eine \moveAction. Lina wählt den Zauber, bezahlt 1 Mana, bestimmt den benachbarten Verbündeten als gültiges Ziel und führt anschließend den Effekt aus. Der Verbündete erhält 4\lifeGain{} zurück.
\end{beispiel}

\section{Angriffszauber}

Enthält der Effekt eines Zaubers eine Angriffsoption, gelten für diese vollständig die Kampfregeln aus \chapterref{chap:combat}. Die Angriffsoption wird erst genutzt, nachdem die Manakosten des Zaubers bezahlt wurden.

\begin{kurzregel}[title=Angriffszauber]
Manakosten bezahlen $\rightarrow$ Angriffsoption nach den normalen Kampfregeln ausführen.
\end{kurzregel}

Bestimmt ein Angriffszauber nach der Angriffsprobe weitere Ziele und überträgt alle gekauften Schadenswürfel und Effekte auf diese, werden die Ziele vor dem Schadenswurf gewählt. Der gemeinsame Schadenspool wird nur einmal gewürfelt und das Ergebnis gegen jedes Ziel angewendet. Schilde und zielabhängige Boni werden für jedes Ziel einzeln berücksichtigt.

\begin{beispiel}[title=Beispiel: Eisblitz]
Lina wirkt \emph{Eisblitz}. Der Zauber kostet 1\manaLoss{} und benötigt eine \standardAction. Sein Effekt ist eine Angriffsoption mit der Angriffsfertigkeit \emph{Zaubern}. Bei dieser Angriffsoption kann 1\star{} für 1\redDmgDie{} oder 2\star{} für 2\redDmgDie{} ausgegeben werden. Für 3\star{} kann das Ziel stattdessen den Statuseffekt \emph{Festgefroren} erhalten.

Lina bezahlt 1 Mana, bestimmt einen gültigen Gegner als Ziel und legt anschließend die Angriffsprobe mit \emph{Zaubern} ab. Sie erzielt 2\star{} und gibt sie für 2\redDmgDie{} aus. Der weitere Ablauf folgt den normalen Regeln für Schaden, Rüstungsbruch, \armor{} und Lebenspunktverlust.
\end{beispiel}

\section{Zauberproben}

Nicht jeder Zauber verlangt eine Probe. Verlangt ein Effekt eine Probe, nennt die Karte Attribute, Fertigkeit und Schwierigkeit. Misslingt die Probe, tritt die Wirkung nicht ein; das ausgegebene Mana bleibt verbraucht.

Eine Angriffsprobe durch eine Angriffsoption ist ebenfalls eine Zauberprobe. Für sie gelten jedoch die Regeln für Angriffe aus \chapterref{chap:combat}.

\section{Ziele und Wirkungsdauer}

Die Zauberkarte bestimmt mögliche Ziele. Typische Zielangaben sind Selbst, Berührung, ein Ziel, mehrere Ziele oder ein Bereich. Ein Ziel muss grundsätzlich wahrnehmbar sein.

Die Reichweite einer Angriffsoption ergibt sich aus ihrer Angriffsfertigkeit, sofern der Zauber nichts anderes bestimmt.

Typische Wirkungsdauern sind sofort, bis zum Ende des Zuges, bis zum Ende der Runde oder eine bestimmte Dauer. Eine Kampfrunde endet, nachdem alle Gegner vollständig aktiviert wurden.

Endet ein Zauber, enden auch alle durch ihn verursachten fortdauernden Boni, Mali und Zustände, sofern die Karte nichts anderes bestimmt.

\section{Ausführliches Beispiel}

Lina möchte im Kampf zwei entfernte Gegner mit \emph{Eisblitz} angreifen.

\begin{enumerate}
  \item \textbf{Zauber wählen:} Lina wählt den verfügbaren Zauber \emph{Eisblitz}. Das Wirken benötigt eine \standardAction.
  \item \textbf{Mana bezahlen:} Sie bezahlt die Grundkosten von 1\manaLoss.
  \item \textbf{Ziel bestimmen:} Da die Angriffsoption \emph{Zaubern} verwendet, darf sie einen benachbarten oder entfernten Gegner als erstes Ziel wählen.
  \item \textbf{Angriffsprobe:} Lina legt eine Angriffsprobe mit \emph{Zaubern} ab und erzielt 3\star.
  \item \textbf{Sterne ausgeben:} Lina gibt 1\star{} für 1\redDmgDie{} und 2\star{} für 2 weitere \redDmgDie{} aus. Ihr Schadenspool besteht aus 3\redDmgDie{}.
  \item \textbf{Weiteres Ziel wählen:} Lina bezahlt zusätzlich 1\manaLoss{} und wählt einen zweiten entfernten Gegner in Reichweite. Beide Ziele sind von allen gekauften Schadenswürfeln betroffen.
  \item \textbf{Angriff abschließen:} Lina würfelt den gemeinsamen Schadenspool einmal. Das Ergebnis wird gegen beide Gegner angewendet; \armor{} und mögliche weitere Effekte werden für jedes Ziel einzeln berücksichtigt.
\end{enumerate}

\section{Weiterführende Informationen}

Die vollständige Zauberliste befindet sich in \appendixref{app:spells}.

\chapter{Was sind Talente?}

Talente verleihen einem Helden besondere Fähigkeiten, die über seine Attribute und Fertigkeiten hinausgehen. Sie machen jeden Helden einzigartig und ermöglichen es ihm, sich im Laufe seiner Abenteuer weiterzuentwickeln.

Es gibt zwei Arten von Talenten:

\begin{itemize}
  \item Rollentalente
  \item allgemeine Talente
\end{itemize}

Während Rollentalente bestimmen, welche Aufgabe ein Held innerhalb seiner Gruppe übernimmt, verleihen allgemeine Talente ihm besondere Fähigkeiten, Eigenschaften oder Kenntnisse.

Die vollständige Übersicht aller Talente befindet sich in Anhang 3 -- Talente.

\section{Was ist ein Talent?}

Ein Talent ist eine besondere Fähigkeit eines Helden.

Talente werden durch Talentkarten dargestellt. Auf jeder Talentkarte steht, welche Wirkung das Talent besitzt und unter welchen Bedingungen es eingesetzt werden kann.

Talente können einem Helden beispielsweise erlauben,

\begin{itemize}
  \item besondere Angriffe auszuführen,
  \item Schaden zu verhindern,
  \item Verbündete zu unterstützen,
  \item Gegner zu behindern,
  \item Zauber zu wirken oder zu verbessern,
  \item außergewöhnliche Fähigkeiten einzusetzen oder
  \item außerhalb eines Kampfes neue Möglichkeiten zu erhalten.
\end{itemize}

Jedes Talent gehört entweder zu den Rollentalenten oder zu den allgemeinen Talenten.

\section{Welche Arten von Talenten gibt es?}

In Heldenpfade gibt es zwei unterschiedliche Talentarten.

Rollentalente bestimmen, welche Aufgabe ein Held innerhalb seiner Gruppe übernimmt. Sie sind vor allem auf den taktischen Teil des Spiels ausgerichtet.

Allgemeine Talente verleihen einem Helden besondere Fähigkeiten, Eigenschaften oder Kenntnisse. Sie beschreiben den Helden selbst und sind nicht an eine bestimmte Rolle gebunden.

Ein Held kann Talente beider Arten miteinander kombinieren.

\section{Was sind Rollentalente?}

Rollentalente bestimmen, auf welche Weise ein Held seiner Gruppe besonders hilft.

Sie sind in vier Rollen unterteilt:

\begin{itemize}
  \item \textbf{Draufgänger} gehen entschlossen in die Offensive. Ihre Talente helfen dabei, viel Schaden zu verursachen und Gegner schnell auszuschalten.
  \item \textbf{Verteidiger} schützen sich selbst und ihre Verbündeten. Sie verhindern Schaden, halten Gegner auf und schaffen Freiräume für die Gruppe.
  \item \textbf{Anführer} stärken ihre Verbündeten und unterstützen die gesamte Gruppe. Sie ermöglichen zusätzliche Handlungen, helfen bei der Heilung oder verbessern die Fähigkeiten anderer Helden.
  \item \textbf{Taktiker} beeinflussen das Kampfgeschehen. Sie behindern Gegner, verändern Positionen und schaffen günstige Situationen für ihre Verbündeten.
\end{itemize}

Eine Rolle ist kein fester Beruf und keine starre Klasse.

Ein Held muss sich nicht dauerhaft für eine einzige Rolle entscheiden. Er kann Rollentalente aus verschiedenen Rollen erlernen und miteinander kombinieren. Dadurch entstehen sehr unterschiedliche Spielstile.

\section{Was sind allgemeine Talente?}

Allgemeine Talente verleihen einem Helden besondere Fähigkeiten, Eigenschaften oder Kenntnisse.

Sie sind nicht an eine bestimmte Rolle gebunden und beschreiben, was einen Helden besonders macht.

Allgemeine Talente können sowohl außerhalb als auch innerhalb eines Kampfes Auswirkungen besitzen. Sie dienen jedoch nicht dazu, eine bestimmte Rolle innerhalb der Gruppe zu erfüllen.

Allgemeine Talente können beispielsweise:

\begin{itemize}
  \item den Umgang mit Magie ermöglichen oder verbessern,
  \item körperliche Fähigkeiten erweitern,
  \item besondere Bewegungsarten verleihen,
  \item außergewöhnliche Sinne ermöglichen,
  \item soziale Fähigkeiten stärken oder
  \item besonderes Wissen und handwerkliche Kenntnisse vermitteln.
\end{itemize}

Dadurch können sich zwei Helden mit derselben Rolle völlig unterschiedlich entwickeln.

\section{Trainierte und untrainierte Talente}

Ein Held kann im Laufe seiner Abenteuer viele Talente erlernen.

Gleichzeitig kann er jedoch nur einen Teil davon aktiv nutzen.

Ein Held kann gleichzeitig bis zu

\begin{itemize}
  \item vier Rollentalente und
  \item vier allgemeine Talente
\end{itemize}

trainiert haben.

Nur trainierte Talente können eingesetzt werden.

Auch nicht trainierte Talente gelten weiterhin als erlernt. Sie können deshalb Voraussetzungen für andere Talente erfüllen.

Wie ein Held neue Talente erlernt, wird in Kapitel 9 -- Wie entsteht ein Held? beschrieben.

\section{Talente neu trainieren}

Im Laufe seiner Abenteuer kann ein Held lernen, andere Talente einzusetzen.

Erhält ein Held die Gelegenheit, neu zu trainieren, darf er beliebig viele seiner trainierten Talente austauschen.

Dabei gelten folgende Regeln:

\begin{itemize}
  \item Es können höchstens vier Rollentalente trainiert sein.
  \item Es können höchstens vier allgemeine Talente trainiert sein.
  \item Rollentalente können nur gegen andere bereits erlernte Rollentalente ausgetauscht werden.
  \item Allgemeine Talente können nur gegen andere bereits erlernte allgemeine Talente ausgetauscht werden.
\end{itemize}

Nach dem Training können die neu gewählten Talente sofort eingesetzt werden.

\section{Wie werden Talente eingesetzt?}

Wie ein Talent eingesetzt wird, steht auf seiner Talentkarte.

Manche Talente wirken dauerhaft. Solange ein Talent trainiert ist, gilt der beschriebene Effekt jederzeit.

Andere Talente werden nur in bestimmten Situationen ausgelöst. Dafür kann zum Beispiel ein Angriff, eine Probe, ein Treffer oder eine andere Handlung notwendig sein.

Manche Talente verlangen außerdem, dass der Held:

\begin{itemize}
  \item einen Mondkristall ausgibt,
  \item die Talentkarte umdreht,
  \item eine bestimmte Handlung verwendet oder
  \item eine andere auf der Talentkarte genannte Bedingung erfüllt.
\end{itemize}

Wird eine Talentkarte umgedreht, ist sie erschöpft. Sie kann erst wieder verwendet werden, nachdem sie bereitgemacht wurde.

\subsection{Mehrere Talente gleichzeitig einsetzen}

Ein Held kann mehrere trainierte Talente besitzen, die in derselben Situation eingesetzt werden können.

Solange die Bedingungen aller Talente erfüllt sind, können ihre Wirkungen miteinander kombiniert werden.

Dabei wird jedes Talent einzeln ausgeführt. Kosten müssen für jedes Talent getrennt bezahlt werden.

Wenn sich zwei Talente widersprechen, gilt die speziellere Regel der Talentkarte.

\section{Der eigene Heldenpfad}

Talente bestimmen den Weg, den ein Held im Laufe seiner Abenteuer einschlägt.

Rollentalente legen fest, welche Aufgabe ein Held innerhalb seiner Gruppe übernimmt. Allgemeine Talente verleihen ihm darüber hinaus besondere Fähigkeiten und Eigenschaften.

Dadurch können zwei Helden mit ähnlichen Attributen, Fertigkeiten und derselben Ausrüstung dennoch vollkommen unterschiedlich sein.

Ein Held kann sich auf eine einzelne Rolle spezialisieren oder Rollentalente aus mehreren Rollen miteinander verbinden. Gleichzeitig kann er durch allgemeine Talente besondere Fähigkeiten entwickeln, die unabhängig von seiner Rolle sind.

Es gibt dabei keinen vorgeschriebenen richtigen Weg. Die Talente zeigen, welche Erfahrungen ein Held gesammelt hat und zu welcher Art von Held er im Laufe seiner Abenteuer geworden ist.
\chapter{Abenteuer erleben}
\label{chap:adventures}

Neben Kämpfen erleben Helden viele weitere Situationen. Sie schleichen an Wachen vorbei, erklimmen steile Felswände, überwinden Hindernisse oder unterstützen sich gegenseitig bei schwierigen Aufgaben.

Dieses Kapitel beschreibt die Regeln für Zusammenarbeit und Rast. Wie genau eine gemeinsame Handlung in der jeweiligen Situation umgesetzt wird, entscheidet die Spielleitung gemeinsam mit der Gruppe.

\section{Zusammenarbeit}

Helden können Herausforderungen entweder gemeinsam bewältigen oder sich gegenseitig unterstützen.

\subsection{Gruppenproben}

Müssen mehrere Helden gemeinsam dieselbe Herausforderung meistern, kann die Spielleitung eine Gruppenprobe verlangen. Sie bestimmt, welche Helden beteiligt sind und welche Proben zur Situation passen.

\begin{kurzregel}[title=Gruppenprobe]
Die Gesamtschwierigkeit entspricht der regulären Schwierigkeit multipliziert mit der Anzahl der beteiligten Helden.

Alle beteiligten Helden legen die von der Spielleitung bestimmten Proben ab und addieren ihre erzielten \star{}. Erreichen oder überschreiten sie gemeinsam die Gesamtschwierigkeit, ist die Gruppenprobe erfolgreich.
\end{kurzregel}

\begin{beispiel}[title=Beispiel: Gemeinsam schleichen]
Vier Helden möchten sich gemeinsam an einer Wache vorbeischleichen. Die Schwierigkeit beträgt 4.

Da vier Helden beteiligt sind, beträgt die Gesamtschwierigkeit 16\star. Alle vier Helden führen eine passende Schleichen-Probe durch. Erzielen sie gemeinsam mindestens 16\star, bleibt die Gruppe unentdeckt.
\end{beispiel}

\subsection{Helfen}

Ein Held kann einen anderen bei einer Probe unterstützen. Dazu beschreibt der Spieler, wie sein Held helfen möchte. Die Spielleitung entscheidet, ob die Unterstützung sinnvoll ist und welche Probe dafür abgelegt wird.

Mehrere Helden können helfen, wenn jeder von ihnen auf eine eigene sinnvolle Weise zur Lösung beiträgt. Ob eine bestimmte Unterstützung möglich ist, entscheidet die Spielleitung.

\begin{kurzregel}[title=Bei einer Probe helfen]
Der Helfer legt zuerst eine passende Probe ab. Für jeden erzielten \star{} erhält der aktive Held einen zusätzlichen \blueDie{} für seine Probe.

Erzielt der Helfer keinen \star{}, erhält der aktive Held keinen Bonus. Weitere Folgen einer misslungenen Unterstützung bestimmt die Spielleitung anhand der Situation.
\end{kurzregel}

\begin{beispiel}[title=Beispiel: Eine Steintür öffnen]
Gadhra hilft Kiera dabei, eine schwere Steintür zu öffnen. Sie beschreibt, wie sie die Tür mit einem Balken anhebt. Die Spielleitung verlangt dafür eine Attributsprobe mit Kraft.

Gadhra erzielt 2\star. Kiera erhält dadurch 2 zusätzliche \blueDie{} für ihre Probe zum Öffnen der Tür.
\end{beispiel}

\section{Rast}

Helden können rasten, wenn sie ausreichend Zeit und Sicherheit dafür haben. Ob und wann eine Rast möglich ist, entscheidet die Spielleitung.

Jeder Held rastet unabhängig von den anderen. Während einige Helden rasten, können andere beispielsweise Wache halten oder anderen Tätigkeiten nachgehen.

Während einer Rast sind nur leichte Tätigkeiten möglich. Kämpfen, Reisen, anstrengende Proben oder vergleichbare Belastungen unterbrechen die Rast. Im Zweifel entscheidet die Spielleitung, ob eine Tätigkeit noch als Rast gilt.

\begin{kurzregel}[title=Eine Stunde rasten]
Für jede vollständig abgeschlossene Stunde Rast erhält ein Held:

\begin{itemize}
  \item Lebenspunkte in Höhe seines KON-Wertes, höchstens bis zu seinem Maximum,
  \item Mana in Höhe seines CHA-Wertes, höchstens bis zu seinem Maximum,
  \item alle erschöpften Karten wieder bereitgemacht.
\end{itemize}
\end{kurzregel}

Wird eine Rast vor Ablauf einer vollständigen Stunde unterbrochen, erhält der Held für diese angefangene Stunde keine Vorteile.

Nach einer vollständig abgeschlossenen Stunde werden alle umgedrehten Karten -- beispielsweise Talente, Zauber und Ausrüstung -- wieder aufgedeckt, sofern eine Regel nichts anderes bestimmt.

\chapter{Wie entsteht ein Held?}

Jeder Held beginnt seine Reise als Abenteurer.

Bevor ein Held sein erstes Abenteuer erlebt, erhält er eine Geschichte, Werte, Fähigkeiten und eine passende Ausrüstung. Diese Entscheidungen bestimmen, wer der Held ist, worin er besonders gut ist und auf welche Weise er Herausforderungen meistert.

Im Laufe seiner Abenteuer sammelt ein Held Erfahrung. Dadurch verbessert er seine Fähigkeiten, erlernt neue Talente und Zauber und entwickelt sich vom Abenteurer über den Helden bis hin zu einer Legende.

Dieses Kapitel beschreibt sowohl die Erschaffung eines neuen Helden als auch seine Weiterentwicklung während einer Kampagne.

\section{Einen Helden erschaffen}

Die Heldenerschaffung erfolgt in mehreren Schritten. Die folgende Checkliste führt durch den gesamten Ablauf.

\begin{kurzregel}[title=Checkliste: Einen Helden erschaffen]
\begin{itemize}[label=$\square$]
  \item Hintergrund und Geschichte ausdenken
  \item Volk wählen
  \item Rolle wählen
  \item Attribute verteilen
  \item Fertigkeiten verteilen
  \item Abgeleitete Werte berechnen
  \item Ein Starttalent wählen
  \item Startzauber wählen
  \item Passende Startausrüstung auswählen
  \item Persönlichkeit und Aussehen ausgestalten
\end{itemize}
\end{kurzregel}

Sind alle Schritte abgeschlossen, ist der Held bereit für sein erstes Abenteuer.

\subsection{Hintergrund und Geschichte ausdenken}

Überlege zunächst, wer dein Held ist und welches Leben er geführt hat, bevor sein erstes Abenteuer beginnt.

Woher kommt er? Was hat er bisher getan? Warum zieht er ins Abenteuer? Welche Ziele verfolgt er, und welche Menschen, Orte oder Ereignisse haben ihn geprägt?

Der Hintergrund muss nicht bis ins letzte Detail feststehen. Einige Antworten können sich erst während des Spiels entwickeln. Er sollte jedoch eine Vorstellung davon vermitteln, wer der Held ist und warum seine späteren Entscheidungen zu ihm passen.

\subsection{Volk wählen}

Wähle, welchem Volk dein Held angehört. Welche Völker in einem Abenteuer vorkommen und welche Bedeutung sie in der Spielwelt besitzen, legen die Spielgruppe und die Spielleitung gemeinsam fest.

Die Wahl des Volkes ist Teil der Geschichte des Helden. Sie kann seine Herkunft, seine Erfahrungen und die Art beeinflussen, wie andere Figuren ihm begegnen.

\subsection{Rolle wählen}

Wähle eine Rolle, die beschreibt, auf welche Weise dein Held seiner Gruppe zu Beginn besonders helfen möchte: Draufgänger, Verteidiger, Anführer oder Taktiker.

Die Rolle ist weder ein Beruf noch eine feste Klasse. Sie dient als Ausgangspunkt für die Entwicklung des Helden. Später kann er Rollentalente aus verschiedenen Rollen erlernen und miteinander kombinieren.

Die Rollen und ihre Schwerpunkte werden in Kapitel 7 -- Was sind Talente? beschrieben.

\subsection{Attribute verteilen}

Zu Beginn der Heldenerschaffung erhält der Spieler acht grüne Attributswürfel, die frei auf die acht Attribute verteilt werden können.

\begin{kurzregel}[title=Attribute verteilen]
Verteile acht grüne Attributswürfel auf die Attribute. Kein Attribut darf zu Beginn mehr als zwei Würfel besitzen. Attribute dürfen auch keinen Würfel besitzen.
\end{kurzregel}

Durch die Verteilung legt der Spieler die natürlichen Stärken und Schwächen seines Helden fest.

\begin{beispiel}[title=Kieras Attribute]
Kiera soll mutig, aufmerksam, magisch begabt und beweglich sein. Sie investiert jeweils zwei Würfel in MUT und INT sowie jeweils einen Würfel in CHA, GEW, KON und KRA.

Damit besitzt sie MUT 2, KLU 0, INT 2, CHA 1, FIN 0, GEW 1, KON 1 und KRA 1. Insgesamt hat sie alle acht Attributswürfel verteilt.
\end{beispiel}

\subsection{Fertigkeiten verteilen}

Zu Beginn der Heldenerschaffung erhält der Spieler neun blaue Fertigkeitswürfel, die frei auf die Fertigkeiten verteilt werden.

\begin{kurzregel}[title=Fertigkeiten verteilen]
Verteile neun blaue Fertigkeitswürfel auf die Fertigkeiten. Keine Fertigkeit darf zu Beginn mehr als zwei Würfel besitzen. Fertigkeiten dürfen auch keinen Würfel besitzen.
\end{kurzregel}

Die Fertigkeiten bestimmen, in welchen Bereichen ein Held bereits Erfahrung gesammelt hat. Eine Übersicht aller Fertigkeiten befindet sich in Kapitel 2 -- Was macht einen Helden aus? Die Regeln für Fertigkeitsproben stehen in Kapitel 3 -- Wie funktionieren Proben?

\subsection{Abgeleitete Werte berechnen}

Nachdem die Attribute verteilt wurden, werden die davon abgeleiteten Werte bestimmt.

\begin{kurzregel}[title=Abgeleitete Werte]
\textbf{Lebenspunkte} $= 14 + 2 \times \mathrm{KON}$

\textbf{Manavorrat} $= 2 \times \mathrm{CHA}$
\end{kurzregel}

Talente und andere Effekte können diese Werte anschließend verändern. Die genauen Regeln für Lebenspunkte und Mana stehen in Kapitel 2 -- Was macht einen Helden aus?

\subsection{Ein Starttalent wählen}

Jeder Held beginnt das Spiel mit einem Starttalent. Als Starttalent darf ein beliebiges Talent gewählt werden, dessen Voraussetzungen der Held beim Spielstart erfüllt.

Das Starttalent gilt als erlernt und trainiert. Die Regeln für Talente sowie die Bedeutung der Rollen stehen in Kapitel 7 -- Was sind Talente?

\subsection{Startzauber wählen}

Hat ein Held das Talent \emph{Magiebegabt} gewählt, kann er bereits zu Beginn Zauber beherrschen.

\begin{kurzregel}[title=Startzauber]
Die Summe der Manakosten aller Startzauber darf den Manavorrat des Helden nicht überschreiten. Kein einzelner Startzauber darf mehr als 4 Mana kosten. Jeder Zauber kann nur einmal erlernt werden.
\end{kurzregel}

Das Talent \emph{Magiebegabt} erhöht den Manavorrat des Helden um 6 Mana. Die Regeln für Zauber stehen in Kapitel 6 -- Wie wirkt Magie?

\subsection{Passende Startausrüstung auswählen}

Jeder Held beginnt das Spiel mit einer Ausrüstung, die zu seiner bisherigen Geschichte passt.

Der Spieler stellt die Startausrüstung gemeinsam mit der Spielleitung zusammen. Sie sollte zu einem neuen Helden auf dem Rang Abenteurer sowie zu seinem sozialen und finanziellen Hintergrund passen. Die Spielleitung entscheidet, welche Gegenstände für den Helden glaubwürdig und angemessen sind.

Ein Jäger könnte beispielsweise einen Bogen und leichte Kleidung besitzen. Eine ehemalige Wache könnte mit Schwert, Schild und Rüstung beginnen. Ein wandernder Heiler trägt vielleicht einen Stab, einen Beutel mit Kräutern und einfache Reisekleidung.

\subsection{Persönlichkeit und Aussehen ausgestalten}

Zum Abschluss erhält der Held die Merkmale, durch die er am Spieltisch lebendig wird. Dazu können unter anderem gehören:

\begin{itemize}
  \item Name und Alter,
  \item Haar- und Augenfarbe,
  \item Körperbau und auffällige Merkmale,
  \item Kleidung und persönlicher Stil,
  \item Eigenheiten, Vorlieben und Abneigungen,
  \item Beziehungen zu anderen Figuren sowie
  \item Hoffnungen, Ängste und persönliche Ziele.
\end{itemize}

Nicht jedes Detail muss bereits vor dem ersten Abenteuer feststehen. Die Figur darf sich während des gemeinsamen Spiels weiterentwickeln.

\section{Beispiel: Kiera erschaffen}

Das folgende Beispiel führt die Heldenerschaffung einmal vollständig durch.

\begin{beispiel}[title=Kiera erschaffen]
\textbf{Hintergrund:} Kiera wuchs als Tochter einer einfachen Handwerkerfamilie auf und arbeitete einige Jahre als Botin. Auf ihren Wegen lernte sie, Gefahren früh zu erkennen, sich unauffällig zu bewegen und notfalls mit einem Speer zu verteidigen. Seit sie erstmals unkontrolliert Magie gewirkt hat, sucht sie nach Antworten über die Herkunft ihrer Kräfte.

\textbf{Volk:} Kiera ist ein Mensch.

\textbf{Rolle:} Sie beginnt als Draufgängerin. Kiera möchte Gegner entschlossen angreifen, ist aber nicht dauerhaft an diese Rolle gebunden.

\textbf{Attribute:} Kiera verteilt ihre acht Attributswürfel auf MUT 2, INT 2, CHA 1, GEW 1, KON 1 und KRA 1. KLU und FIN bleiben auf 0.

\textbf{Fertigkeiten:} Sie verteilt ihre neun Fertigkeitswürfel auf Schlagen 2, Zaubern 2, Schleichen 2, Turnen 1, Suchen 1 und Überreden 1. Alle übrigen Fertigkeiten bleiben auf 0.

\textbf{Abgeleitete Werte:} Mit KON 1 besitzt Kiera $14 + 2 \times 1 = 16$ Lebenspunkte. Mit CHA 1 besitzt sie zunächst 2 Mana.

\textbf{Starttalent:} Kiera wählt \emph{Magiebegabt}. Das Talent erhöht ihren Manavorrat um 6, sodass sie mit 8 Mana beginnt.

\textbf{Startzauber:} Kiera erlernt \emph{Weißer Speer} für 1 Mana, \emph{Zielscheibe} für 2 Mana, \emph{Geisterschritt} für 3 Mana und \emph{Herzschlag spüren} für 2 Mana. Die Summe ihrer Manakosten entspricht genau ihrem Manavorrat von 8.

\textbf{Startausrüstung:} Gemeinsam mit der Spielleitung wählt Kiera einen Speer und eine Knochenrüstung. Diese Ausrüstung passt zu ihrer bisherigen Geschichte und bereitet sie auf ihr erstes Abenteuer vor.

\textbf{Persönlichkeit und Aussehen:} Kiera ist neugierig, direkt und ungeduldig. Sie trägt ihr dunkles Haar meist zusammengebunden, bevorzugt praktische Kleidung und bewahrt eine alte Messingbrosche ihrer Mutter auf. Sie möchte lernen, ihre Magie zu beherrschen, und fürchtet zugleich, durch sie Menschen zu gefährden, die ihr nahestehen.
\end{beispiel}

\section{Einen Helden weiterentwickeln}

Während seiner Abenteuer sammelt ein Held Erfahrung. Mit dieser Erfahrung kann er seine Fähigkeiten verbessern, neue Talente und Zauber erlernen und sich zu einem mächtigeren Helden entwickeln.

\subsection{Erfahrungspunkte erhalten}

Die Spielleitung vergibt Erfahrungspunkte für bestandene Abenteuer, besondere Leistungen oder andere bedeutende Entwicklungen.

Wie genau Erfahrungspunkte vergeben werden, entscheidet jede Spielleitung. Für eine durchschnittliche Kampagne wird empfohlen:

\begin{itemize}
  \item \textbf{Gruppen-EP:} Alle Helden erhalten gleichzeitig die gleiche Menge an EP.
  \item \textbf{EP-Menge:} Etwa 1 bis 2 EP pro Spielabend.
\end{itemize}

\subsection{Heldenränge}

Mit zunehmender Erfahrung steigt ein Held im Rang auf. Es gibt drei Heldenränge: Abenteurer, Held und Legende.

Für den Heldenrang zählt die gesamte Erfahrung, die ein Held im Laufe seiner Abenteuer erhalten hat. Bereits ausgegebene Erfahrungspunkte werden weiterhin mitgezählt.

Jeder Heldenrang erhöht die maximal möglichen Attributs- und Fertigkeitswerte.

\begin{hptable}{L{2.5cm} C{3.0cm} C{3.1cm} C{3.1cm}}
\textbf{Heldenrang} & \textbf{Ab erhaltenen EP} & \textbf{Max. Attributswert} & \textbf{Max. Fertigkeitswert} \\
Abenteurer & 0 & 3 & 3 \\
Held & 75 & 4 & 4 \\
Legende & 150 & 5 & 5 \\
\end{hptable}

Ein Held muss den entsprechenden Heldenrang erreicht haben, bevor er ein Attribut oder eine Fertigkeit über den bisherigen Höchstwert hinaus steigern darf.

\subsection{Erfahrungspunkte ausgeben}

Ein Held kann seine gesammelten Erfahrungspunkte verwenden, um sich weiterzuentwickeln.

\begin{hptable}{Y L{4.5cm}}
\textbf{Verbesserung} & \textbf{Kosten} \\
Fertigkeit steigern & Neuer Fertigkeitswert $\times$ 1 EP \\
Attribut steigern & Neuer Attributswert $\times$ 3 EP \\
Talent erlernen & siehe unten \\
Zauber erlernen & Manakosten $\times$ 2 EP \\
\end{hptable}

Die Kosten zum Erlernen neuer Talente richten sich danach, wie viele Talente der Held bereits erlernt hat.

\begin{hptable}{Y L{4.5cm}}
\textbf{Erlerntes Talent} & \textbf{Kosten} \\
Starttalent & 0 EP \\
2. Talent & 5 EP \\
3. Talent & 10 EP \\
4. Talent & 15 EP \\
Jedes weitere Talent & 15 EP \\
\end{hptable}

Neue Talente gelten sofort als trainiert. Sind dadurch bereits vier Talente derselben Talentart trainiert, entscheidet der Spieler, welches seiner bisher trainierten Talente dieser Talentart stattdessen nicht mehr trainiert ist.

\appendix
\chapter{Ausrüstungsgegenstände}
\label{app:equipment}

Dieser Anhang enthält die derzeit verfügbaren Ausrüstungsgegenstände. Die Bedeutung der verwendeten Symbole ist in \appendixref{app:symbol-reference} erläutert.

\section{Waffen}

Waffen können eine oder mehrere Angriffsoptionen besitzen. Jede Angriffsoption nennt die verwendete Angriffsfertigkeit, die verfügbaren Kaufoptionen und gegebenenfalls besondere Voraussetzungen.

\subsection{Allgemeine Kampfeffekte}

Kampfeffekte sind standardisierte Regeln des Kampfsystems. Sie können durch Angriffsoptionen, Talente, Zauber, Gegnerfähigkeiten oder andere Spieleffekte ausgelöst werden. Jeder Kampfeffekt ist an dieser Stelle definiert und wird an anderer Stelle lediglich referenziert.

\begin{hptable}{L{3.0cm} Y}
\textbf{Kampfeffekt} & \textbf{Wirkung} \\
\hline
Rüstungsbruch & Erhöht den Rüstungsbruch des Angriffs um den angegebenen Wert. \\
Defensiv & Der Angreifer erhält bis zum Beginn seines nächsten Zuges den bei der Kaufoption angegebenen Bonus auf Ausweichen oder die angegebenen Schild-Symbole. \\
Spalten & Nicht für das ursprüngliche Ziel ausgegebene Sterne dürfen für Schadenswürfel gegen einen weiteren Gegner in Reichweite verwendet werden. \\
Zu Fall bringen & Das Ziel erhält den Statuseffekt \emph{Zu Boden}. \\
Wegstoßen & Das Ziel wird ein Feld direkt vom Angreifer weg bewegt. War es zuvor benachbart, gilt es anschließend als entfernt. \\
Betäuben & Das Ziel setzt seinen nächsten Zug vollständig aus. \\
\end{hptable}

\begin{hptable}{L{2.6cm} C{1.7cm} L{2.7cm} L{3.2cm} Y}
\textbf{Name} & \textbf{Plätze} & \textbf{Angriffsfertigkeit} & \textbf{Schaden} & \textbf{Effekt} \\
Schwert & 1 Hand & Schlagen & \star{} $\rightarrow$ \orangeDmgDie & -- \\
Speer & 2 Hände & Schlagen & \star{} $\rightarrow$ \orangeDmgDie & \armorPierce{} 2 \\
 & & Schießen & \star{} $\rightarrow$ \orangeDmgDie\orangeDmgDie & \exhaust \\
Zahnschwert & 2 Hände & Schlagen & \star{} $\rightarrow$ \orangeDmgDie\orangeDmgDie & -- \\
\end{hptable}

\begin{todo}[title=Offen: Handlungsart besonderer Ausrüstungseffekte]
Für Ausrüstungseffekte wie die Heilung der Knochenrüstung ist noch festzulegen, ob ihre Nutzung grundsätzlich eine freie Aktion ist oder ob die Handlungsart jeweils auf dem Gegenstand angegeben wird.
\end{todo}

\section{Rüstungen und Schilde}

Rüstungen und Schilde können Schild-Symbole verleihen, eine Ausweichprobe ermöglichen oder verbessern und besondere Effekte besitzen.

\begin{hptable}{L{2.8cm} C{1.7cm} C{2.0cm} L{4.1cm} Y}
\textbf{Name} & \textbf{Platz} & \textbf{Schild-Symbole} & \textbf{Ausweichen} & \textbf{Effekt} \\
Kettenrüstung & Körper & \armor\armor & \heroDie{} + INT + GEW & -- \\
Knochenrüstung & Körper & \armor & \heroDie{} + INT + GEW & \exhaust{} $\rightarrow$ 8\,\lifeGain \\
Kleiner Deckel & 1 Hand & -- & +\blueDie{} auf \evade & -- \\
Großer Deckel & 1 Hand & \armor & -- & -- \\
\end{hptable}

\section{Sonstige Gegenstände}

Derzeit sind keine sonstigen Ausrüstungsgegenstände aufgeführt.

\chapter{Zaubersprüche}
\label{app:spells}

Dieser Anhang enthält alle Zauber, die Helden im Laufe ihrer Abenteuer erlernen können.

Die Zauber sind nach ihren \textbf{Manakosten} sortiert. Jeder Zauber nennt seine Schlüsselworte, seine Zauberzeit und seine Wirkung. Erfordert ein Zauber eine Zauberprobe, ist dies in seiner Effektbeschreibung angegeben.

Bei Angriffszaubern werden \star{} wie bei Waffen gemäß den aufgeführten Zeilen ausgegeben. Das Schlüsselwort \emph{Elementar} bezeichnet die Zauberschule; \emph{Feuer}, \emph{Eis}, \emph{Wind} und \emph{Erde} bezeichnen das verwendete Element.

\section{Zauber für 1 Mana}

\begin{hptable}{L{2.7cm} L{3.0cm} C{1.0cm} L{2.1cm} Y}
\textbf{Zauber} & \textbf{Schlüsselworte} & \textbf{Mana} & \textbf{Zauberzeit} & \textbf{Effekt} \\
\textbf{Flammenstrahl} & Elementar, Feuer & 1 & \standardAction{} & \emph{Zaubern}: \star{} $\rightarrow$ \redDmgDie{}\newline 2\star{} $\rightarrow$ \redDmgDie{}\redDmgDie{}\newline \star{} $\rightarrow$ Das Ziel erhält \emph{Brennend 1}.\newline 1\manaLoss{} $\rightarrow$ Wähle ein weiteres, bisher nicht gewähltes Ziel in Reichweite. Es ist von allen gekauften Schadenswürfeln und Effekten betroffen. Diese Option darf mehrfach gewählt werden. \\
\textbf{Eisblitz} & Elementar, Eis & 1 & \standardAction{} & \emph{Zaubern}: \star{} $\rightarrow$ \redDmgDie{}\newline 2\star{} $\rightarrow$ \redDmgDie{}\redDmgDie{}\newline 3\star{} $\rightarrow$ Das Ziel erhält \emph{Festgefroren}.\newline 1\manaLoss{} $\rightarrow$ Wähle ein weiteres, bisher nicht gewähltes Ziel in Reichweite. Es ist von allen gekauften Schadenswürfeln und Effekten betroffen. Diese Option darf mehrfach gewählt werden. \\
\textbf{Windklinge} & Elementar, Wind & 1 & \standardAction{} & \emph{Zaubern}: \star{} $\rightarrow$ \redDmgDie{}\newline 2\star{} $\rightarrow$ \redDmgDie{}\redDmgDie{}\newline \star{} $\rightarrow$ \effectPush{}\newline 1\manaLoss{} $\rightarrow$ Wähle ein weiteres, bisher nicht gewähltes Ziel in Reichweite. Es ist von allen gekauften Schadenswürfeln und Effekten betroffen. Diese Option darf mehrfach gewählt werden. \\
\textbf{Steinschlag} & Elementar, Erde & 1 & \standardAction{} & \emph{Zaubern}: \star{} $\rightarrow$ \redDmgDie{}\newline 2\star{} $\rightarrow$ \redDmgDie{}\redDmgDie{}\newline 2\star{} $\rightarrow$ \effectKnockdown{}\newline 1\manaLoss{} $\rightarrow$ Wähle ein weiteres, bisher nicht gewähltes Ziel in Reichweite. Es ist von allen gekauften Schadenswürfeln und Effekten betroffen. Diese Option darf mehrfach gewählt werden. \\
\textbf{Weißer Speer} & Beschwörung & 1 & \freeAction{} & Besitzt du einen Speer, wird dieser bereitgemacht. Andernfalls erscheint ein Speer in deinen Händen, sofern deine Hände-Plätze frei sind. \\
\textbf{Heilender Funke} & Leben & 1 & \moveAction{} & Ein benachbarter Verbündeter erhält 4\lifeGain{} zurück. \\
\end{hptable}

\section{Zauber für 2 Mana}

\begin{hptable}{L{2.7cm} L{3.0cm} C{1.0cm} L{2.1cm} Y}
\textbf{Zauber} & \textbf{Schlüsselworte} & \textbf{Mana} & \textbf{Zauberzeit} & \textbf{Effekt} \\
\textbf{Herzschlag spüren} & Leben, Erkenntnis & 2 & \standardAction{} & Finde alle Lebewesen innerhalb von 10 Metern. \\
\textbf{Zielscheibe} & Erkenntnis & 2 & \standardAction{} & Du erhältst bis zum Ende des Kampfes +1\blueDie{} auf Angriffsproben gegen einen Gegner deiner Wahl. \\
\end{hptable}

\section{Zauber für 3 Mana}

\begin{hptable}{L{2.7cm} L{3.0cm} C{1.0cm} L{2.1cm} Y}
\textbf{Zauber} & \textbf{Schlüsselworte} & \textbf{Mana} & \textbf{Zauberzeit} & \textbf{Effekt} \\
\textbf{Schneesturm} & Elementar, Eis & 3 & \standardAction{} & Wähle einen sichtbaren Bereich von bis zu 3 $\times$ 5 Metern. Durch den Bereich kann nicht gesehen werden. Angriffe und Zauber, die Sicht auf ihr Ziel benötigen, können nicht durch den Bereich hindurch ausgeführt werden. Betritt eine Figur den Bereich oder bewegt sie sich darin, muss sie einmal pro Bewegung eine schwierige Probe mit \emph{Turnen} ablegen. Misslingt die Probe, geht sie zu Boden. Beginnt eine brennende Figur ihren Zug im Schneesturm, endet \emph{Brennend}. Zu Beginn der nächsten Runde endet der Schneesturm, sofern der Zaubernde nicht 1\manaLoss{} bezahlt. Dies darf zu Beginn jeder folgenden Runde wiederholt werden. \\
\textbf{Geisterschritt} & Bewegung, Jenseits & 3 & \standardAction{} & \emph{Zaubern}: Erschaffe ein Portal. Es verschwindet, sobald eine Anzahl verbündeter Wesen in Höhe deiner erzielten \star{} es durchschritten hat oder dieselbe Anzahl Minuten vergangen ist. Das Portal kann von beiden Seiten durchquert werden. \\
\textbf{Geisterblick} & Erkenntnis, Jenseits & 3 & \standardAction{} & Du siehst für 5 Minuten Geheimtüren und Geheimgänge. \\
\end{hptable}

\section{Zauber für 4 Mana}

\begin{hptable}{L{2.7cm} L{3.0cm} C{1.0cm} L{2.1cm} Y}
\textbf{Zauber} & \textbf{Schlüsselworte} & \textbf{Mana} & \textbf{Zauberzeit} & \textbf{Effekt} \\
\textbf{Freundschaft} & Einfluss & 4 & \standardAction{} & Eine Kreatur betrachtet dich für 5 Minuten als Freund. Im Kampf setzt ein Gegner für 1 Runde aus. \\
\end{hptable}

\begin{todo}
Widerstand, gültige Ziele und Grenzen des Zaubers \emph{Freundschaft} festlegen.
\end{todo}

\section{Zauber für 5 Mana}

\begin{hptable}{L{2.7cm} L{3.0cm} C{1.0cm} L{2.1cm} Y}
\textbf{Zauber} & \textbf{Schlüsselworte} & \textbf{Mana} & \textbf{Zauberzeit} & \textbf{Effekt} \\
\textbf{Feuerball} & Elementar, Feuer & 5 & \standardAction{} & Wähle einen sichtbaren Punkt in bis zu 20 Metern Entfernung. Würfle einmal \redDmgDie{}\redDmgDie{}\redDmgDie{}\redDmgDie{}\redDmgDie{}. Alle Wesen im Umkreis von 3 Metern um diesen Punkt erleiden das gewürfelte Ergebnis und erhalten \emph{Brennend 2}. Freund und Feind sind gleichermaßen betroffen. Diesem Schaden kann nicht ausgewichen werden. Brennbare Objekte im Wirkungsbereich fangen Feuer. \\
\end{hptable}

\section{Zauber für 6 Mana}

\begin{hptable}{L{2.7cm} L{3.0cm} C{1.0cm} L{2.1cm} Y}
\textbf{Zauber} & \textbf{Schlüsselworte} & \textbf{Mana} & \textbf{Zauberzeit} & \textbf{Effekt} \\
\multicolumn{5}{c}{\emph{Noch keine Zauber eingetragen.}} \\
\end{hptable}

\section{Zauber für 7 Mana}

\begin{hptable}{L{2.7cm} L{3.0cm} C{1.0cm} L{2.1cm} Y}
\textbf{Zauber} & \textbf{Schlüsselworte} & \textbf{Mana} & \textbf{Zauberzeit} & \textbf{Effekt} \\
\multicolumn{5}{c}{\emph{Noch keine Zauber eingetragen.}} \\
\end{hptable}

\section{Zauber für 8 oder mehr Mana}

\begin{hptable}{L{2.7cm} L{3.0cm} C{1.0cm} L{2.1cm} Y}
\textbf{Zauber} & \textbf{Schlüsselworte} & \textbf{Mana} & \textbf{Zauberzeit} & \textbf{Effekt} \\
\multicolumn{5}{c}{\emph{Noch keine Zauber eingetragen.}} \\
\end{hptable}
\chapter{Talente}
\label{app:talents}

Dieser Anhang enthält alle allgemeinen Talente und Rollentalente, die Helden erlernen können.

\begin{kurzregel}[title=Talente einsetzen]
Nennt ein nicht dauerhaftes Talent weder eine Aktion noch eine Reaktion, wird es im eigenen Zug als \freeAction{} eingesetzt. Eine Reaktion ist eine \freeAction{} außerhalb des eigenen Zuges und darf nur unmittelbar nach ihrem fest definierten Auslöser eingesetzt werden. Alle angegebenen Voraussetzungen und Kosten müssen vollständig erfüllt sein.
\end{kurzregel}

\section{Allgemeine Talente}

\subsection{Magietalente}

\begin{hptable}{L{2.8cm} L{2.7cm} L{2.1cm} Y}
\textbf{Name} & \textbf{Voraussetzung} & \textbf{Einsatz} & \textbf{Effekt} \\
\hline
Magiebegabt & CHA 1+ & Dauerhaft & Du darfst Zauber erlernen und wirken. Dein Manavorrat erhöht sich um 6. \\
Manavorrat & Magiebegabt & Dauerhaft & Dein Manavorrat erhöht sich um 2. \\
Schnelles Zaubern & Magiebegabt & Dauerhaft & Wähle beim Trainieren ein Schlüsselwort. Die Zauberzeit aller Zauber mit diesem Schlüsselwort wird im Kampf um eine Aktionsstufe beschleunigt: \standardAction{} wird zu \moveAction{}, \moveAction{} wird zu \freeAction{}. Eine \freeAction{} kann nicht weiter beschleunigt werden. Außerhalb eines Kampfes wirkst du diese Zauber 30\,\% schneller. \\
Sparsames Zaubern & Magiebegabt & Dauerhaft & Wähle beim Trainieren ein Schlüsselwort. Für alle Zauber mit diesem Schlüsselwort musst du beim Wirken 1 Mana weniger ausgeben, mindestens jedoch 1 Mana. Trainierst du das Talent neu, darfst du ein neues Schlüsselwort wählen. \\
\end{hptable}

\subsection{Körper und Bewegung}

\begin{hptable}{L{2.8cm} L{2.7cm} L{2.1cm} Y}
\textbf{Name} & \textbf{Voraussetzung} & \textbf{Einsatz} & \textbf{Effekt} \\
\hline
Akrobat & GEW 2+ & Dauerhaft & Vorteil auf Proben mit \emph{Turnen} beim Springen, Balancieren oder kontrollierten Fallen. \\
Sprinter & GEW 2+ & \exhaust & Erhalte sofort eine zusätzliche \moveAction{}, die du unmittelbar nutzen darfst. \\
Standfest & KON 2+ & Reaktion + \exhaust & Würdest du den Statuseffekt \emph{Zu Boden} erhalten, darfst du als Reaktion die Talentkarte erschöpfen und verhinderst den Statuseffekt. \\
Schnell auf den Beinen & GEW 1+ & 1 \crystal & Stehe sofort auf, ohne dafür eine \moveAction{} auszugeben. \\
\end{hptable}

\subsection{Wahrnehmung und Natur}

\begin{hptable}{L{2.8cm} L{2.7cm} L{2.1cm} Y}
\textbf{Name} & \textbf{Voraussetzung} & \textbf{Einsatz} & \textbf{Effekt} \\
\hline
Spurenleser & INT 1+ & Dauerhaft & Vorteil auf Proben mit \emph{Natur}, um Spuren zu finden, zu verfolgen oder daraus Informationen abzuleiten. \\
Orientierungssinn & INT 1+ & Dauerhaft & Du kennst unter gewöhnlichen Bedingungen stets die Himmelsrichtungen und kannst einen selbst zurückgelegten Weg wiederfinden. Ist dennoch eine Probe zur Orientierung nötig, hast du Vorteil. \\
Tierversteher & CHA 1+ & Dauerhaft & Vorteil auf Proben mit \emph{Natur}, um Tiere zu beruhigen, ihr Verhalten zu deuten oder mit ihnen umzugehen. Ihre grundsätzliche Stimmung und offensichtlichen Bedürfnisse erkennst du ohne Probe. \\
\end{hptable}

\subsection{Soziales}

\begin{hptable}{L{2.8cm} L{2.7cm} L{2.1cm} Y}
\textbf{Name} & \textbf{Voraussetzung} & \textbf{Einsatz} & \textbf{Effekt} \\
\hline
Menschenkenntnis & KLU 2+ & Dauerhaft & Vorteil auf Proben mit \emph{Herausfinden}, um Absichten, Gefühle oder Lügen einer Person zu erkennen. Eine offensichtliche Stimmung erkennst du ohne Probe. \\
Wortgewandt & CHA 2+ & Dauerhaft & Vorteil auf Proben mit \emph{Überreden}, wenn dein Gegenüber bereit ist, dir zuzuhören und ausreichend Zeit für ein Gespräch besteht. \\
Einschüchternd & MUT 2+ & Dauerhaft & Vorteil auf Proben mit \emph{Überreden}, wenn du eine glaubwürdige Drohung aussprichst oder deine Überlegenheit demonstrierst. \\
Bühnenprofi & CHA 2+ & Dauerhaft & Wähle beim Trainieren eine Form des Auftretens, etwa Musik, Schauspiel, Tanz oder Erzählkunst. Vorteil auf passende Proben mit \emph{Vorführen}. \\
Mitreißend & Bühnenprofi & \exhaust & Nachdem dir eine Probe mit \emph{Vorführen} gelungen ist, darfst du die Talentkarte erschöpfen. Bis zu CHA Personen, die deinen Auftritt miterlebt haben, erhalten Vorteil auf ihre nächste Probe mit einer Fertigkeit der Kategorie \emph{Miteinander} in dieser Szene. \\
\end{hptable}

\subsection{Handwerk, Heilkunde und Wissen}

\begin{hptable}{L{2.8cm} L{2.7cm} L{2.1cm} Y}
\textbf{Name} & \textbf{Voraussetzung} & \textbf{Einsatz} & \textbf{Effekt} \\
\hline
Handwerksmeister & FIN 2+ & Dauerhaft & Wähle beim Trainieren ein Handwerk. Vorteil auf passende Proben mit \emph{Basteln}, um entsprechende Gegenstände herzustellen, umzubauen oder zu reparieren. \\
Improvisateur & Handwerksmeister & \standardAction{} + \exhaust & Stelle aus greifbaren Materialien einen einfachen provisorischen Gegenstand her. Er ermöglicht eine klar bestimmte Verwendung und wird anschließend unbrauchbar. Die Spielleitung entscheidet, was mit den vorhandenen Materialien möglich ist. \\
Ersthelfer & KLU 2+ & \standardAction{} + \exhaust & Ein benachbarter kampfunfähiger Held heilt 1 Lebenspunkt und ist dadurch nicht länger kampfunfähig. \\
Erfahrener Heiler & Ersthelfer & 1 \crystal & Wenn du einer Kreatur Lebenspunkte zurückgibst, darfst du 1 \crystal{} ausgeben. Die Kreatur heilt 1 zusätzlichen Lebenspunkt. \\
Gelehrter & KLU 2+ & Dauerhaft & Wähle beim Trainieren ein Wissensgebiet. Vorteil auf passende Proben mit \emph{Herausfinden}. \\
Geistesblitz & Gelehrter & \exhaust & Nachdem eine Probe mit \emph{Herausfinden} misslungen ist, darfst du die Talentkarte erschöpfen. Du erhältst von der Spielleitung einen wahren, hilfreichen Hinweis. Der Hinweis löst das Problem nicht unmittelbar, zeigt aber einen möglichen nächsten Ansatz. \\
\end{hptable}

\clearpage
\section{Rollentalente}

\subsection{Draufgänger}

\begin{hptable}{L{2.8cm} L{2.7cm} L{2.1cm} Y}
\textbf{Name} & \textbf{Voraussetzung} & \textbf{Einsatz} & \textbf{Effekt} \\
\hline
Perfekter Treffer & -- & 1 \crystal & Nachdem alle Schadenspools eines deiner Angriffe geworfen wurden, aber bevor der Schaden ermittelt wird, darfst du beliebig viele der dabei geworfenen Schadenswürfel auf ihre höchste Seite drehen. \\
Mehrfachangriff & Perfekter Treffer & 1 \crystal & Nachdem alle Sterne ausgegeben und Schadenswürfel verbessert wurden, aber bevor ein Schadenspool verteilt oder geworfen wird, wähle einen zweiten Gegner in Angriffsreichweite und verdopple den vollständigen Schadenspool. Ohne Spalten wird je eine Hälfte des Pools einem der beiden Gegner zugeteilt. Besitzt der Angriff Spalten, dürfen anschließend alle Würfel des verdoppelten Schadenspools nach den Regeln für Spalten verteilt werden. Schild-Symbole und sonstige Abwehrregeln werden für jeden Gegner separat angewendet. \\
Schwachstelle & Perfekter Treffer & 1 \crystal & Nachdem das Ziel eines Angriffs bestimmt wurde, kannst du für diesen Angriff alle Schild-Symbole dieses Gegners ignorieren. \\
Blutrausch & Perfekter Treffer & Dauerhaft & Hast du höchstens die Hälfte deiner Lebenspunkte, werden nach dem Ausgeben aller Sterne und dem normalen Verbessern der Schadenswürfel alle deine Schadenswürfel um eine weitere Stufe verbessert: \orangeDmgDie{} $\rightarrow$ \redDmgDie{}, \redDmgDie{} $\rightarrow$ \blackDmgDie{}, \blackDmgDie{} bleibt \blackDmgDie{}. Erst danach werden Schadenspools durch Mehrfachangriff oder andere Effekte verändert und gegebenenfalls verteilt. \\
Zielen & Perfekter Treffer & \standardAction & Verzichte auf deinen Angriff. Bei deinem nächsten Angriff in diesem Kampf zählt jede leere Seite deiner \blueDie{} und \greenDie{} als ein \star. \\
\end{hptable}

\subsection{Verteidiger}

\begin{hptable}{L{2.8cm} L{2.7cm} L{2.1cm} Y}
\textbf{Name} & \textbf{Voraussetzung} & \textbf{Einsatz} & \textbf{Effekt} \\
\hline
Beschützen & -- & 1 \crystal & Wähle einen benachbarten Gegner. Dieser muss dich mit seiner nächsten Aktion angreifen. Kann er das nicht, verfällt seine nächste Aktion. \\
Bollwerk & Beschützen & Reaktion & Sobald ein benachbarter Gegner versucht, deine Nähe zu verlassen, darfst du als Reaktion vor Abschluss seiner Bewegung einen Gelegenheitsangriff mit einer aktuell verfügbaren Angriffsoption für Schlagen oder Zaubern ausführen. Stammt die Angriffsoption von einem Zauber, musst du dessen Manakosten bezahlen und alle weiteren Voraussetzungen erfüllen; die Zauberzeit wird durch die Reaktion ersetzt. Der erste erzielte \star{} beendet die Bewegung und steht für den Angriff nicht mehr zur Verfügung. Mit allen weiteren \star{} darfst du nach den normalen Regeln Sterne gemäß der gewählten Angriffsoption ausgeben und Schadenswürfel verbessern. Erzielst du keinen \star, setzt der Gegner seine Bewegung fort. \\
Schildwall & Beschützen & Dauerhaft & Benachbarte Verbündete können bei Angriffen gegen sie deine Schild-Symbole anstelle ihrer eigenen verwenden. \\
Verteidigen & Beschützen & \exhaust & Verdopple deine Schild-Symbole bis zum Beginn deines nächsten Zugs. \\
Unbeugsam & Beschützen & Reaktion + \exhaust & Würdest du auf 0 oder weniger Lebenspunkte sinken, darfst du als Reaktion die Talentkarte erschöpfen. Du sinkst stattdessen auf $2 \times \mathrm{KON}$ Lebenspunkte. \\
\end{hptable}

\subsection{Anführer}

\begin{hptable}{L{2.8cm} L{2.7cm} L{2.1cm} Y}
\textbf{Name} & \textbf{Voraussetzung} & \textbf{Einsatz} & \textbf{Effekt} \\
\hline
Heilende Worte & -- & \standardAction & Heile einen benachbarten Helden um dein CHA an Lebenspunkten. Gibst du währenddessen 1 \crystal{} aus, verdoppelt sich die Heilung auf $2 \times \mathrm{CHA}$. \\
Ermutigen & Heilende Worte & Reaktion + 1 \crystal & Nachdem ein Verbündeter seine Probe geworfen hat, darfst du als Reaktion 1 \crystal{} ausgeben. Der Verbündete darf alle \blueDie{} dieser Probe neu würfeln. \\
Inspiration & Heilende Worte & Dauerhaft & Gibst du einen \crystal{} aus, darfst du ihn anschließend an einen Verbündeten deiner Wahl weitergeben. Ein Mondkristall darf insgesamt nur einmal weitergegeben werden; wurde er bereits durch Inspiration oder einen anderen Effekt weitergegeben, kann er nicht erneut weitergegeben werden. \\
Erneuern & Heilende Worte & 3 \crystal{} + \exhaust & Mache eine erschöpfte Talentkarte eines Verbündeten bereit. \\
Aufmuntern & Heilende Worte & \exhaust & Richte einen kampfunfähigen Helden mit vollen Lebenspunkten wieder auf. \\
\end{hptable}

\subsection{Taktiker}

\begin{hptable}{L{2.8cm} L{2.7cm} L{2.1cm} Y}
\textbf{Name} & \textbf{Voraussetzung} & \textbf{Einsatz} & \textbf{Effekt} \\
\hline
Taktieren & -- & 1 \crystal & Sieh dir die nächsten drei KI-Einträge eines Gegners an. Entscheide, zu welchem Eintrag die KI springt. \\
Zu Fall bringen & Taktieren & 1 \crystal & Ein benachbarter Gegner erhält den Kampfeffekt \effectKnockdown{} \emph{Zu Fall bringen} und damit den Statuseffekt \emph{Zu Boden}. \\
Umpositionieren & Taktieren & 1 \crystal & Wähle bis zu zwei dir benachbarte Helden. Sie erhalten jeweils eine kostenlose \moveAction{}, die sie sofort nutzen dürfen. \\
Voraussicht & Taktieren & Dauerhaft & Du darfst bei deinen Proben nach dem Wurf beliebig viele \star{} als \moon{} werten. \\
Meisterplan & Taktieren & \exhaust & Alle Helden führen sofort einen zusätzlichen Zug aus. Anschließend führen die Gegner zwei vollständige Gegnerphasen aus. \\
\end{hptable}

\begin{todo}
Die Funktionsweise von \emph{Taktieren} muss mit dem finalen Gegner- und KI-System abgeglichen werden.
\end{todo}
\chapter{Effekte und Zustände}
\label{app:status-effects}

Dieser Anhang enthält standardisierte Kampfeffekte und Statuseffekte. Sie können durch Waffen, Zauber, Talente, Gegnerfähigkeiten oder andere Spielregeln ausgelöst werden und gelten unabhängig von ihrer Quelle immer gleich.

\section{Allgemeine Kampfeffekte}

Kampfeffekte sind standardisierte Regeln des Kampfsystems. Sie können durch Angriffsoptionen, Talente, Zauber, Gegnerfähigkeiten oder andere Spieleffekte ausgelöst werden. Jeder Kampfeffekt ist an dieser Stelle definiert und wird an anderer Stelle lediglich referenziert.

\begin{hptable}{L{4.0cm} Y}
\textbf{Kampfeffekt} & \textbf{Wirkung} \\
\hline
\effectArmorBreak{} \textbf{Rüstungsbruch} & Ignoriert bei der Schadensabwicklung die angegebene Anzahl an Schild-Symbolen jedes betroffenen Ziels. \\
\effectDefensive{} \textbf{Defensiv} & Der Angreifer erhält bis zum Beginn seines nächsten Zuges den angegebenen Bonus auf Ausweichen oder die angegebenen Schild-Symbole. \\
\effectCleave{} \textbf{Spalten} & Vor dem Schadenswurf dürfen die Schadenswürfel des Angriffs auf beliebig viele Gegner in Reichweite verteilt werden. Die einem Gegner zugeteilten Würfel bilden einen eigenen Schadenspool. Rüstungsbruch wird bei jedem dieser Schadenspools berücksichtigt. Alle weiteren ausgelösten Kampfeffekte gelten für jeden betroffenen Gegner, der durch seinen Schadenspool mindestens 1 Lebenspunkt verliert. Wird der Schadenspool zuvor durch Mehrfachangriff verdoppelt, werden anschließend alle Würfel des verdoppelten Pools verteilt. \\
\effectKnockdown{} \textbf{Zu Fall bringen} & Das Ziel erhält den Statuseffekt \emph{Zu Boden}. \\
\effectPush{} \textbf{Wegstoßen} & Das Ziel wird ein Feld direkt vom Angreifer weg bewegt. War es zuvor benachbart, gilt es anschließend als entfernt. \\
\effectStun{} \textbf{Betäuben} & Das Ziel setzt seinen nächsten Zug vollständig aus. \\
\end{hptable}

\section{Statuseffekte}

Statuseffekte beschreiben besondere Zustände, die eine Figur vorübergehend beeinflussen.

Ein Statuseffekt gilt so lange, bis die bei ihm beschriebene Bedingung erfüllt ist oder ein anderer Effekt ihn beendet.

\begin{hptable}{L{3.0cm} Y L{4.0cm}}
\textbf{Statuseffekt} & \textbf{Wirkung} & \textbf{Ende} \\
\hline
Brennend X & Erhält eine Figur \emph{Brennend X}, erhält sie den angegebenen Wert. Ist sie bereits brennend, erhöht sich ihr vorhandener Wert stattdessen um X. Wird der Schaden eines Feuerzaubers gegen eine brennende Figur ermittelt, erhöht sich der Schaden einmalig um ihren Brennend-Wert. Anschließend wird der Wert durch neu erhaltenes Brennend erhöht. & Die Figur verwendet eine \standardAction{}, um Brennend unabhängig vom Wert vollständig zu entfernen, oder sie wird von einem Zauber mit dem Schlüsselwort \emph{Eis} betroffen. \\
Festgefroren & Die Figur kann sich nicht bewegen oder bewegt werden. Ein festgefrorener Held hat Nachteil auf eigene Angriffsproben und auf Ausweichproben. Gegen einen festgefrorenen Gegner haben Helden Vorteil auf Angriffsproben und auf Ausweichproben gegen dessen Angriffe. & Am Ende des nächsten Zuges der Figur oder sobald sie von einem Zauber mit dem Schlüsselwort \emph{Feuer} betroffen wird. \\
Kampfunfähig & Ein kampfunfähiger Held kann nicht handeln. & Sobald der Held geheilt wird und mindestens 1\lifeGain{} besitzt. Anschließend erhält er den Statuseffekt \emph{Zu Boden}. \\
Zu Boden & Helden haben Vorteil auf Angriffsproben gegen zu Boden gegangene Gegner und auf Ausweichproben gegen deren Angriffe. Ein zu Boden gegangener Held hat Nachteil auf eigene Angriffsproben und auf seine Ausweichprobe (\evade). Die Figur muss ihre nächste \moveAction{} verwenden, um aufzustehen. & Sobald die Figur aufgestanden ist. \\
\end{hptable}

\begin{todo}
\textbf{Weitere Statuseffekte ergänzen}

Dieser Anhang wird mit weiteren Zaubern, Talenten und Gegnern erweitert. Mögliche spätere Statuseffekte sind unter anderem Vergiften, Verlangsamen, Furcht und Versteinern.
\end{todo}
\chapter{Beispielhelden}

% Wiederverwendbare Darstellung fuer Beispielhelden

\definecolor{hpHero}{HTML}{F2EEE6}
\definecolor{hpHeroLine}{HTML}{8A6A49}
\definecolor{hpHeroInner}{HTML}{FBF9F4}

% Geschlossenes Profil fuer einen Beispielhelden.
\newtcolorbox{heroprofile}[2][]{
  hpbox,
  breakable=false,
  title={#2},
  colback=hpHero,
  colframe=hpHeroLine,
  boxed title style={colback=hpHeroLine!22},
  fonttitle=\HPHeadingFont\large,
  #1
}

% Acht gleich breite Attributfelder direkt unter dem Namen.
\newenvironment{heroattributes}
  {\begingroup\HPUIFont\renewcommand{\arraystretch}{1.15}\setlength{\tabcolsep}{2pt}%
   \tabularx{\linewidth}{@{}*{8}{>{\centering\arraybackslash}X}@{}}}
  {\endtabularx\endgroup}

% Zweispaltiger Hauptbereich: links Fertigkeiten, rechts weitere Heldenwerte.
\newenvironment{herocolumns}
  {\par\medskip\noindent
   \begin{minipage}[t]{0.485\linewidth}\vspace{0pt}}
  {\end{minipage}\par}

\newcommand{\herocolumnbreak}{%
  \end{minipage}\hfill
  \begin{minipage}[t]{0.485\linewidth}\vspace{0pt}%
}

% Frei stapelbare Informationskarte innerhalb eines Heldenprofils.
\newtcolorbox{herosectionbox}[1][]{
  enhanced,
  colback=hpHeroInner,
  colframe=hpLine,
  boxrule=0.55pt,
  arc=1.5mm,
  left=3mm,
  right=3mm,
  top=2.5mm,
  bottom=2.5mm,
  before skip=0pt,
  after skip=3mm,
  fontupper=\HPUIFont,
  #1
}

% Kompakte Fertigkeitstabelle. Die beiden rechten Spalten zeigen
% Grundwuerfel (Held + Attribute) und zusaetzliche Fertigkeitswuerfel.
\newenvironment{heroskills}
  {\begingroup\HPUIFont\renewcommand{\arraystretch}{1.12}\setlength{\tabcolsep}{3pt}%
   \tabularx{\linewidth}{@{}X >{\centering\arraybackslash}p{2.45cm} >{\centering\arraybackslash}p{1.3cm}@{}}}
  {\endtabularx\endgroup}

\newcommand{\heroskillheader}{%
  & \heroDie{}\,\greenDie{} & \blueDie{} \\
  \hline
}


Dieser Anhang enthält vollständig ausgearbeitete Beispielhelden. Sie können direkt für ein Abenteuer verwendet oder als Vorlage für eigene Helden genutzt werden.

\begin{heroprofile}{Kiera}
  \begin{heroattributes}
    \textbf{MUT} & \textbf{KLU} & \textbf{INT} & \textbf{CHA} & \textbf{FIN} & \textbf{GEW} & \textbf{KON} & \textbf{KRA} \\
    2~\greenDie & -- & 2~\greenDie & -- & -- & 1~\greenDie & 2~\greenDie & 1~\greenDie
  \end{heroattributes}

  \begin{herocolumns}
    \begin{herosectionbox}
      \begin{heroskills}
        \heroskillheader
        Schlagen & \heroDie{} + 3~\greenDie & 2~\blueDie \\
        Schießen & \heroDie{} + 1~\greenDie & -- \\
        Zaubern & \heroDie{} + 2~\greenDie & 2~\blueDie \\
        Turnen & \heroDie{} + 2~\greenDie & 1~\blueDie \\
        Schleichen & \heroDie{} + 3~\greenDie & 2~\blueDie \\
        Suchen & \heroDie{} + 2~\greenDie & 1~\blueDie \\
        Nachdenken & \heroDie{} + 2~\greenDie & 1~\blueDie \\
        Natur kennen & \heroDie{} + 2~\greenDie & 1~\blueDie \\
        Basteln & \heroDie & -- \\
        Überzeugen & \heroDie{} + 2~\greenDie & 1~\blueDie \\
        Heilen & \heroDie{} + 2~\greenDie & 1~\blueDie \\
        Vorführen & \heroDie & --
      \end{heroskills}
    \end{herosectionbox}

    \herocolumnbreak

    \begin{herosectionbox}
      \begin{tabularx}{\linewidth}{@{}>{\bfseries}l X@{}}
        Lebenspunkte & 18~\lifeGain \\
        Mana & 6~\manaGain
      \end{tabularx}
    \end{herosectionbox}

    \begin{herosectionbox}
      \textbf{Talent}\par\smallskip
      Magiebegabt
    \end{herosectionbox}

    \begin{herosectionbox}
      \textbf{Ausrüstung}\par\smallskip
      Speer\\
      Knochenrüstung
    \end{herosectionbox}
  \end{herocolumns}
\end{heroprofile}

\chapter{Gegner-Almanach}

% Wiederverwendbare Darstellung fuer Gegnerprofile

\definecolor{hpEnemy}{HTML}{F2EEE6}
\definecolor{hpEnemyLine}{HTML}{8A6A49}
\definecolor{hpEnemyInner}{HTML}{FBF9F4}

% Ein Gegnerprofil bildet einen geschlossenen Almanach-Eintrag.
\newtcolorbox{enemyprofile}[2][]{
  hpbox,
  title={#2},
  colback=hpEnemy,
  colframe=hpEnemyLine,
  boxed title style={colback=hpEnemyLine!22},
  fonttitle=\HPHeadingFont\large,
  #1
}

% Kompakte Teilbereiche innerhalb eines Gegnerprofils.
\newtcolorbox{enemysectionbox}[1]{
  enhanced,
  colback=hpEnemyInner,
  colframe=hpLine,
  boxrule=0.55pt,
  arc=1.5mm,
  left=3mm,
  right=3mm,
  top=2.5mm,
  bottom=2.5mm,
  before skip=0pt,
  after skip=0pt,
  title={#1},
  fonttitle=\HPUIFont\bfseries,
  fontupper=\HPUIFont,
  coltitle=hpBrown,
  attach title to upper={\quad}
}

% Werteband direkt unter dem Gegnernamen.
\newenvironment{enemystats}
  {\begingroup\HPUIFont\renewcommand{\arraystretch}{1.15}\setlength{\tabcolsep}{4pt}%
   \tabularx{\linewidth}{@{}>{\bfseries}l X >{\bfseries}l X >{\bfseries}l X@{}}}
  {\endtabularx\endgroup}

% Zweispaltiger Inhaltsbereich fuer Angriffe, KI und Sonderregeln.
\newenvironment{enemygrid}
  {\begin{tcbraster}[
      raster columns=2,
      raster column skip=4mm,
      raster row skip=3mm,
      raster valign=top
    ]}
  {\end{tcbraster}}

\newcommand{\enemyrole}[1]{\textbf{Rolle:} #1}
\newcommand{\enemydesign}[1]{\textbf{Designziel:} #1}


Der Gegner-Almanach versammelt die spielrelevanten Werte, Angriffe und Verhaltensregeln der Gegner. Jeder Eintrag ist als kompaktes Profil aufgebaut: Das Werteband zeigt die wichtigsten Kennzahlen, darunter stehen Angriffe und Verhalten direkt nebeneinander.

\section{Höhlenspinne}

\begin{enemyprofile}{Höhlenspinne}
  \begin{enemystats}
    Lebenspunkte & 12~\lifeLoss & Schilde & 1~\armor & Rolle & Kontrollgegner \\
  \end{enemystats}

  \smallskip
  Eine lauernde Jägerin, die ihre Beute zunächst mit klebrigen Netzen schwächt und anschließend aus nächster Nähe zubeißt.

  \begin{enemygrid}
    \begin{enemysectionbox}{Angriffe}
      \textbf{Netzspucken} \hfill \rangedAttack

      5~\orangeDmgDie{} Netzwürfel

      Jeder nicht negierte Erfolg erzeugt einen Netzmarker auf dem Ziel.

      \medskip
      \textbf{Biss} \hfill \meleeAttack

      3~\redDmgDie{} Schadenswürfel
    \end{enemysectionbox}

    \begin{enemysectionbox}{Verhalten}
      Solange mindestens ein Held keinen Netzmarker besitzt, setzt die Höhlenspinne \textbf{Netzspucken} ein.

      Besitzen alle Helden mindestens einen Netzmarker, führt sie einen \textbf{Biss} aus.
    \end{enemysectionbox}

    \begin{enemysectionbox}{Netzmarker}
      Jeder Netzmarker verhindert den nächsten Angriffserfolg oder Ausweicherfolg des betroffenen Helden.

      Anschließend wird der Marker entfernt.
    \end{enemysectionbox}

    \begin{enemysectionbox}{Taktische Rolle}
      \enemyrole{Kontrolle und Vorbereitung}

      \medskip
      \enemydesign{Die Höhlenspinne verursacht zunächst wenig direkten Schaden. Ihre Netze schwächen jedoch gezielt die nächsten Angriffe oder Ausweichversuche und erhöhen so den Druck nachfolgender Aktionen.}
    \end{enemysectionbox}
  \end{enemygrid}
\end{enemyprofile}

\section{Baumdrache}

\begin{enemyprofile}{Baumdrache}
  \begin{enemystats}
    Lebenspunkte & 50~\lifeLoss & Schilde & 3~\armor & Rolle & Bossgegner \\
  \end{enemystats}

  \smallskip
  Ein mächtiger Gegner mit einer vorhersehbaren, aber abwechslungsreichen Angriffsfolge. Wer seine Muster erkennt, kann sich gezielt auf den nächsten Schlag vorbereiten.

  \begin{enemygrid}
    \begin{enemysectionbox}{Angriffe}
      \textbf{Feueratem} \hfill \rangedAttack

      5~\orangeDmgDie{} gegen alle Helden

      \medskip
      \textbf{Biss} \hfill \meleeAttack

      3~\redDmgDie{}

      \medskip
      \textbf{Schwanzschlag} \hfill \meleeAttack

      3~\orangeDmgDie{}; das Ziel wird weggeschleudert

      \medskip
      \textbf{Krallenhieb} \hfill \meleeAttack

      4~\orangeDmgDie{}
    \end{enemysectionbox}

    \begin{enemysectionbox}{Angriffsfolge}
      Der Baumdrache führt pro Runde genau eine Aktion aus:

      \begin{enumerate}[leftmargin=1.5em,itemsep=0.15\baselineskip]
        \item Feueratem
        \item Biss
        \item Schwanzschlag
        \item Krallenhieb
        \item Biss
        \item Schwanzschlag
        \item Krallenhieb
      \end{enumerate}

      Danach beginnt die Folge erneut.
    \end{enemysectionbox}

    \begin{enemysectionbox}{Prolog-Regel}
      Im Prolog endet der Kampf nicht mit dem Tod des Baumdrachen.

      Sobald er insgesamt \textbf{30 Schaden} erlitten hat, flieht er.
    \end{enemysectionbox}

    \begin{enemysectionbox}{Taktische Rolle}
      \enemyrole{Boss mit lernbarer Rotation}

      \medskip
      \enemydesign{Der Baumdrache verbindet Flächenschaden, hohen Einzelschaden und Positionskontrolle. Seine feste Reihenfolge macht die Bedrohung planbar und belohnt vorausschauendes Handeln.}
    \end{enemysectionbox}
  \end{enemygrid}
\end{enemyprofile}

\chapter{Symbolübersicht}
\label{app:symbol-reference}

Dieser Anhang erklärt die im Regelwerk verwendeten Symbole. Ihre Bedeutung bleibt unabhängig von ihrer konkreten grafischen Gestaltung gleich.

\section{Würfel}

\begin{hptable}{C{2.2cm} Y}
\textbf{Symbol} & \textbf{Bedeutung} \\
\heroDie & Heldenwürfel. Jede Probe eines Helden enthält genau einen Heldenwürfel. \\
\greenDie & Grüner Attributswürfel. Er steht für die natürlichen Eigenschaften eines Helden. \\
\blueDie & Blauer Fertigkeitswürfel. Er steht für Übung, Wissen und Erfahrung in einer Fertigkeit. \\
\orangeDmgDie & Oranger Schadenswürfel für leichten Schaden. \\
\redDmgDie & Roter Schadenswürfel für mittleren Schaden. \\
\blackDmgDie & Schwarzer Schadenswürfel für schweren Schaden. \\
\end{hptable}

\section{Würfelergebnisse und Ressourcen}

\begin{hptable}{C{2.2cm} Y}
\textbf{Symbol} & \textbf{Bedeutung} \\
\star & Stern. Sterne zeigen Erfolge bei Proben und Angriffen. \\
\moon & Mond. Jeder gewürfelte Mond erzeugt sofort einen Mondkristall. \\
\crystal & Mondkristall. Gib ihn aus, um alle Würfel einer Farbe neu zu würfeln oder einen anderen Effekt auszulösen. Ein Held kann beliebig viele Mondkristalle besitzen. Wann ungenutzte Mondkristalle verfallen, bestimmt die Spielleitung. \\
\end{hptable}

\section{Leben und Mana}

\begin{hptable}{C{2.2cm} Y}
\textbf{Symbol} & \textbf{Bedeutung} \\
\lifeLoss & Lebenspunkte verlieren. Die angegebene Zahl wird von den aktuellen Lebenspunkten abgezogen; Lebenspunkte sinken nie unter 0. \\
\lifeGain & Lebenspunkte erhalten. Die angegebene Zahl wird zu den aktuellen Lebenspunkten addiert, höchstens bis zum Maximum. \\
\manaLoss & Mana bezahlen oder verlieren. Die angegebene Zahl wird vom aktuellen Mana abgezogen. \\
\manaGain & Mana erhalten. Die angegebene Zahl wird zum aktuellen Mana addiert, höchstens bis zum Manavorrat. \\
\end{hptable}

\section{Aktionen}

\begin{hptable}{C{2.2cm} Y}
\textbf{Symbol} & \textbf{Bedeutung} \\
\standardAction & Standardaktion. Sie wird für Angriffe, Zauber und andere größere Handlungen verwendet. \\
\moveAction & Bewegungsaktion. Sie wird für Bewegung oder das Neuordnen von Ausrüstung verwendet. \\
\freeAction & Freie Aktion. Sie steht für kurze Handlungen, die kaum Zeit kosten. \\
\end{hptable}

\section{Angriffsfertigkeiten}

\begin{hptable}{C{2.2cm} Y}
\textbf{Symbol} & \textbf{Bedeutung} \\
\meleeAttack & Schlagen. Kennzeichnet Angriffe, die mit der Fertigkeit \textbf{Schlagen} ausgeführt werden. \\
\rangedAttack & Schießen. Kennzeichnet Angriffe, die mit der Fertigkeit \textbf{Schießen} ausgeführt werden. \\
\magicAttack & Zaubern. Kennzeichnet Angriffe, die mit der Fertigkeit \textbf{Zaubern} ausgeführt werden. \\
\end{hptable}

\section{Kampf und Karten}

\begin{hptable}{C{2.2cm} Y}
\textbf{Symbol} & \textbf{Bedeutung} \\
\armor & Schild-Symbol. Jedes Schild-Symbol verhindert 1 Schaden. Schild-Symbole können durch Rüstungen, Schilde, Talente und andere Effekte gewährt werden. \\
\armorPierce & Rüstungsdurchdringung. Die angegebene Anzahl an Schild-Symbolen des Ziels wird ignoriert. \\
\evade & Ausweichen. Das Symbol kennzeichnet Ausweichproben oder Effekte, mit denen Schadenswürfel entfernt werden können. \\
\exhaust & Erschöpfen. Drehe die Karte um. Solange sie erschöpft ist, können ihre Angaben und Effekte nicht genutzt werden. \\
\end{hptable}

\end{document}
